\documentclass[11pt]{amsart}
% Traducao para o portugues do artigo (main.tex). As secoes estao em sec_pt/, com os mesmos rotulos; os
% numeros seguem a versao em ingles (ponto decimal). Compilar: pdflatex main_pt && bibtex main_pt && pdflatex main_pt (x2).
\usepackage[T1]{fontenc}
\usepackage[utf8]{inputenc}
\usepackage[brazilian]{babel}
\usepackage{lmodern}
\usepackage{amsmath,amssymb,amsthm,mathtools,thmtools}
\usepackage[margin=1.15in]{geometry}
\usepackage{enumitem}
\usepackage{booktabs}
\usepackage{array}
\usepackage[dvipsnames]{xcolor}
\usepackage[colorlinks=true,linkcolor=MidnightBlue,citecolor=MidnightBlue,urlcolor=MidnightBlue]{hyperref}
\usepackage[brazilian,capitalise,nameinlink]{cleveref}

\def\versaopt{}
\renewcommand{\datename}{\textit{Data}}
\makeatletter\def\emailaddrname{{\itshape E-mail}}\makeatother
% ---- theorem environments (thmtools, so that cleveref names each kind correctly)
% main_pt.tex (the Portuguese translation) defines \versaopt before reading this file.
\ifdefined\versaopt
\declaretheorem[name=Teorema,refname={Teorema,Teoremas},numberwithin=section]{theorem}
\declaretheorem[name=Proposição,refname={Proposição,Proposições},sibling=theorem]{proposition}
\declaretheorem[name=Lema,refname={Lema,Lemas},sibling=theorem]{lemma}
\declaretheorem[name=Corolário,refname={Corolário,Corolários},sibling=theorem]{corollary}
\declaretheorem[name=Definição,refname={Definição,Definições},sibling=theorem,style=definition]{definition}
\declaretheorem[name=Observação,refname={Observação,Observações},sibling=theorem,style=remark]{remark}
\declaretheorem[name=Teorema Principal,numbered=no]{mainthm}
\declaretheorem[name=Teorema,numbered=no]{introthm}
\else
\declaretheorem[name=Theorem,numberwithin=section]{theorem}
\declaretheorem[name=Proposition,sibling=theorem]{proposition}
\declaretheorem[name=Lemma,sibling=theorem]{lemma}
\declaretheorem[name=Corollary,sibling=theorem]{corollary}
\declaretheorem[name=Definition,sibling=theorem,style=definition]{definition}
\declaretheorem[name=Remark,sibling=theorem,style=remark]{remark}
\declaretheorem[name=Main Theorem,numbered=no]{mainthm}
\declaretheorem[name=Theorem,numbered=no]{introthm}
\fi

% ---- drafting aids (remove before submission)
\newcommand{\todo}[1]{{\color{BrickRed}[\textbf{TODO:} #1]}}
\newcommand{\src}[1]{{\color{Gray}\footnotesize[\texttt{#1}]}}

% ---- notation
\newcommand{\R}{\mathbb{R}}
\newcommand{\C}{\mathbb{C}}
\newcommand{\Z}{\mathbb{Z}}
\renewcommand{\Re}{\operatorname{Re}}
\renewcommand{\Im}{\operatorname{Im}}
\newcommand{\gbar}{\bar g}
\newcommand{\phibar}{\bar\phi}
\newcommand{\omegabar}{\bar\omega}
\newcommand{\Mhat}{\widehat M}
\newcommand{\Lop}{L}             % linearized RT family
\newcommand{\Cop}{C}             % linearized constraints
\newcommand{\Kop}{K}             % constraint propagation operator
\newcommand{\Lie}{\mathcal{L}}
\newcommand{\sphys}{s_{\ast}}    % the physical root
\newcommand{\sK}{s_{K}}
\newcommand{\sB}{s_{B}}
\newcommand{\Dom}{\mathcal{D}}
\newcommand{\Yp}{Y_{+}}
\newcommand{\mult}{\operatorname{m}}

% DOI links in the bibliography (amsplain prints the note field)
\newcommand{\doi}[1]{\href{https://doi.org/#1}{\texttt{doi:#1}}}

\hypersetup{pdftitle={Exatamente um modo instável: uma contagem espectral assistida por computador para o espaço-tempo
crítico de Choptuik sob perturbações esfericamente simétricas},pdfauthor={Edivaldo Amaral Gon\c{c}alves},
pdfkeywords={espaço-tempo de Choptuik, colapso crítico, modos instáveis, prova assistida por computador, teorema de
Rouché, método de Newton-Kantorovich}}

\title[Um modo instável para o espaço-tempo de Choptuik]{Exatamente um modo instável: uma contagem espectral
assistida por computador para o espaço-tempo crítico de Choptuik sob perturbações esfericamente simétricas}
\author{Edivaldo Amaral Gon\c{c}alves}
\address{Instituto Federal de Mato Grosso (IFMT), Cuiab\'a, Mato Grosso, Brasil}
\email{edivaldo.goncalves@ifmt.edu.br}
\urladdr{https://orcid.org/0009-0007-9117-9833}
\thanks{Instituto Federal de Mato Grosso (IFMT), Cuiab\'a, Mato Grosso, Brasil.\newline
E-mail: \texttt{edivaldo.goncalves@ifmt.edu.br}.\quad
ORCID: \href{https://orcid.org/0009-0007-9117-9833}{0009-0007-9117-9833}.\newline
Tradução para o português da versão~2 do preprint em inglês, \emph{Exactly one unstable mode: a computer-assisted
spectral count for Choptuik's critical spacetime under spherically symmetric perturbations}
(DOI \href{https://doi.org/10.5281/zenodo.23251929}{\nolinkurl{10.5281/zenodo.23251929}}, que reúne todas as versões). Em caso de divergência, vale o texto em
inglês. Os números seguem a convenção da versão em inglês, com ponto decimal.}
\date{Preprint, versão 2, outubro de 2026}

\begin{document}

\begin{abstract}
O espaço-tempo de Choptuik, a solução discretamente autossimilar no limiar de formação de buraco negro no
colapso de um campo escalar sem massa, deve ter exatamente um modo instável, cuja taxa de crescimento $\lambda$
determina o expoente crítico $\gamma=1/\lambda\approx0.374$. A existência desse espaço-tempo foi provada por
Reiterer e Trubowitz com um argumento assistido por computador. Provamos, também com assistência computacional, que na
realização deles a solução tem exatamente um modo instável, módulo gauge, sob perturbações lineares esfericamente
simétricas que são suaves no tempo e analíticas na variável radial até o cone de luz e através dele. O
modo é real e simples, com $\lambda\in[2.674040,\,2.674115]$, de modo que
$1/\lambda\in[0.373955,\,0.373966]$. Mais precisamente, por classe de Floquet, o operador linearizado tem exatamente quatro
raízes com parte real positiva: um modo de gauge, que translada o instante da singularidade; duas raízes cujos
autovetores violam os vínculos; e o modo físico. A contagem usa uma versão para operadores do teorema de Rouch\'e
com uma referência de dimensão finita, uma deflação de posto dois das raízes de gauge, localizações rigorosas por
Newton--Kantorovich e cotas em aritmética racional e de bolas. A passagem do problema geométrico ao operador
se apoia num resultado de completude global do gauge e numa reconstrução explícita das perturbações. Os
certificados, um programa de verificação e os dados estão disponíveis publicamente.
\end{abstract}

\maketitle
\setcounter{tocdepth}{1}
\tableofcontents

\section{Introdução}\label{sec:intro}

\subsection{Colapso crítico}

Em 1993, Choptuik estudou numericamente o colapso gravitacional de um campo escalar sem massa esfericamente simétrico
e encontrou, no limiar entre a dispersão e a formação de buraco negro, uma solução universal com três características
marcantes~\cite{Choptuik1993}: ela é aproximada, em todas as famílias a um parâmetro de dados iniciais que foram evoluídas, quando o
parâmetro é ajustado ao limiar; ela é discretamente autossimilar, reproduzindo-se em escalas cada vez menores com um período de eco
$\Delta\approx3.44$; e, perto do limiar, a massa do buraco negro escala como $(p-p_*)^\gamma$ com um
expoente universal $\gamma\approx0.37$. Esses \emph{fenômenos críticos} foram desde então encontrados em muitos modelos de matéria, e são
revisados em~\cite{GundlachMartinGarcia2007}. Para o sistema de Einstein--campo escalar esfericamente simétrico, Christodoulou provou que
dados suficientemente pequenos se dispersam e dados suficientemente grandes formam superfícies aprisionadas~\cite{Christodoulou1986,
Christodoulou1991,Christodoulou1993}.

A explicação padrão é dinâmica~\cite{KoikeHaraAdachi1995,Gundlach1997}. A solução crítica é um ponto
fixo (no caso discretamente autossimilar, uma órbita periódica) da evolução no tempo logarítmico, com uma variedade
estável de codimensão um, que é o próprio limiar, e um \emph{único} modo instável, que cresce como
$e^{\lambda T}$. Uma solução quase crítica fica perto da crítica até que o modo instável cresça até um
tamanho fixo, e a análise dimensional dá então $\gamma=1/\lambda$, com uma modulação periódica no caso discretamente
autossimilar. Gundlach calculou numericamente as perturbações da solução de Choptuik e encontrou exatamente um modo que cresce,
com $\gamma=0.374\pm0.001$~\cite{Gundlach1997}, e Mart\'in-Garc\'ia e Gundlach encontraram numericamente que todas as
perturbações não esféricas decaem~\cite{MartinGarciaGundlach1999}. Um ingrediente necessário desse quadro é que a
solução crítica tenha exatamente um modo instável; é esse enunciado espectral que provamos, para perturbações
esfericamente simétricas.

\subsection{Existência, e a questão espectral em aberto}

A existência do espaço-tempo de Choptuik foi provada por Reiterer e Trubowitz~\cite{RT2019}, com um argumento
assistido por computador na tradição da prova de Lanford das conjecturas de Feigenbaum~\cite{Lanford1982}. Eles constroem uma
solução analítica real e discretamente autossimilar numa vizinhança do cone de luz passado da singularidade, como
ponto fixo de uma contração perto de dados numéricos de alta precisão, e determinam a sua constante de eco e um
segundo invariante com $80$ dígitos. Sobre as perturbações, eles escrevem: ``We do not study perturbations about the
Choptuik spacetime, that are relevant for many of the (conjectured) properties of the solution. Our paper can be
the starting point for a rigorous investigation of this kind. This would be interesting, because current
numerical results seem to be inconclusive''~\cite[\S1]{RT2019} (``Não estudamos perturbações do espaço-tempo de
Choptuik, que são relevantes para muitas das propriedades (conjecturadas) da solução. O nosso artigo pode ser o ponto
de partida para uma investigação rigorosa desse tipo. Isso seria interessante, porque os resultados numéricos atuais
parecem inconclusivos'').

Este artigo faz uma investigação desse tipo para perturbações esfericamente simétricas.

\subsection{O resultado}

Trabalhamos com a solução de~\cite{RT2019}, nas suas coordenadas $(\tau,\xi)$ e variáveis $\omega$, nas quais a
autossimilaridade é a translação $\tau\mapsto\tau+2\pi$. As perturbações lineares do tipo de Floquet,
$e^{s\tau}$ vezes um perfil periódico, são regidas por uma família de operadores $L(s)$, a linearização do
sistema de Reiterer--Trubowitz, junto com uma aplicação de vínculos (\emph{constraints}) linearizada $C(s)$. O expoente $s$ corresponde à
taxa de crescimento $\lambda=2\pi s/K$ no tempo logarítmico usual, onde $e^{-K}$ é o fator pelo qual a
autossimilaridade reescala os comprimentos.

\begin{introthm}[versão informal do Teorema Principal, \cref{sec:main}]
Sob perturbações lineares esfericamente simétricas, suaves no tempo e analíticas na variável radial até o cone de luz
e através dele, e módulo gauge (calibre), o espaço-tempo de Choptuik tem exatamente um modo instável. O modo é real e algebricamente simples, e a sua taxa de crescimento satisfaz
\[
  \lambda\in[2.674040,\ 2.674115],\qquad 1/\lambda\in[0.373955,\ 0.373966].
\]
Mais precisamente, por classe de Floquet, o operador $L$ tem exatamente quatro raízes no semiplano $\Re s>0$, cada uma simples
e com um representante real: um modo de gauge, que é a translação do instante da singularidade; duas raízes cujos autovetores
violam os vínculos; e o modo físico instável. No eixo imaginário, $L$ tem apenas a raiz $s=0$, a
translação da fase. Se o cone de luz da singularidade é mantido fixo, o modo de gauge deixa de ser removido pelo
grupo menor de transformações de gauge, e a contagem passa a ser dois; a classe adicional corresponde a transladar
o instante da singularidade, e não a uma segunda instabilidade.
\end{introthm}

O valor de $1/\lambda$ concorda com o expoente crítico medido e calculado; esta é uma conferência externa, já que
a relação $\gamma=1/\lambda$ diz respeito à evolução não linear. O teorema é espectral: ele localiza as raízes de
$L$ e identifica as físicas, mas não diz nada sobre o semigrupo da evolução linearizada nem sobre a dinâmica não
linear. Ele trata apenas de perturbações esfericamente simétricas, na classe de regularidade enunciada acima. Essas
limitações são discutidas na \cref{sec:discussion}.

\subsection{Estratégia}

A prova tem três partes.

\emph{Do problema geométrico ao operador} (\cref{sec:physical}). Um modo físico é uma solução das equações
de Einstein--campo escalar linearizadas do tipo de Floquet, e só a sua classe módulo transformações de gauge tem
significado. Mostramos que todo modo desses com $\Re s>0$ pode ser posto na forma de~\cite{RT2019} por uma transformação
de gauge, e então corresponde a um elemento de $\ker L(s)\cap\ker C(s)$; que, reciprocamente, todo elemento desses
vem de um modo físico; e que os únicos modos de gauge com $\Re s>0$ são os múltiplos de um modo explícito em
$s=\mu$, a translação infinitesimal do ponto singular. A completude do gauge é global no domínio
de~\cite{RT2019} e usa que o cone de luz está contido nele. A multiplicidade física de uma raiz é então uma questão
de álgebra linear no seu espaço de raízes.

\emph{Contagem} (\cref{sec:counting}). As raízes de $L$ com parte real grande são excluídas por uma cota da
imagem numérica (\emph{numerical range}), calculada num espaço de funções no qual o fundo (\emph{background}, a solução
em torno da qual se lineariza) age ponto a ponto. Na região limitada
restante, as raízes são contadas com uma versão para operadores do teorema de Rouch\'e, devida a Gohberg e
Sigal~\cite{GohbergSigal1971}: após uma deflação (\emph{deflation}) de posto dois, que afasta as duas raízes de gauge, $L$ é comparado ao longo da
fronteira com uma referência bloco-diagonal cujos zeros são os autovalores de uma matriz $14016\times14016$, e estes
são contados a partir de uma decomposição de Schur com resíduo certificado.

\emph{Identificação} (\cref{sec:roots}). Cada raiz é localizada rigorosamente numa bola certificada (\emph{enclosure}), e mostra-se que é
algebricamente simples, por um argumento de Newton--Kantorovich aplicado a um sistema aumentado (\emph{bordered system}):
a equação de autovalor acrescida de uma condição de normalização, com o autovalor como incógnita. Mostra-se que os autovetores em duas delas violam os vínculos, por cotas inferiores
explícitas. Para a raiz restante, os vínculos não podem ser conferidos num autovetor aproximado; em vez disso, uma
identidade $KC=ML$ propaga os vínculos, e um certificado à parte mostra que o operador de propagação dos vínculos
$K$ é injetivo ali.

Todos os cálculos são feitos em aritmética racional exata, em aritmética de bolas (\emph{ball arithmetic}), ou como
cotas a posteriori sobre cálculos em ponto flutuante. Na aritmética de bolas, cada número real é representado por um
centro $m$ e um raio $r$, isto é, pelo intervalo $[m-r,m+r]$, e cada operação devolve uma bola que contém o resultado
exato. Eles levam alguns dias em computadores comuns.

\subsection{O que há de novo}

As ferramentas são clássicas; o que é novo é a sua combinação para este operador, e os resultados. O artigo transforma em
teorema um quadro espectral que só era conhecido numericamente: ele dá uma contagem rigorosa dos modos instáveis do
espaço-tempo de Choptuik; uma localização rigorosa da
taxa de crescimento física; a classificação de todas as raízes com parte real positiva; e a passagem do problema
geométrico ao operador, incluindo a completude do gauge. Os certificados, um programa de verificação que
os refaz, os dados e um guia de reprodução estão disponíveis publicamente (\cref{sec:computation-data}).

\subsection{Organização}

A \cref{sec:setting} recorda a formulação de~\cite{RT2019} e define $L$, $C$ e $K$. A \cref{sec:main} enuncia o
Teorema Principal. A \cref{sec:physical} relaciona os modos físicos ao operador, a \cref{sec:counting} conta as raízes e
a \cref{sec:roots} as identifica e completa a prova. A \cref{sec:computation} descreve os cálculos, a sua
verificação e o uso de assistentes de IA, e a \cref{sec:discussion} discute as limitações e o problema não
esférico. Os apêndices contêm as provas dos lemas estruturais, fórmulas explícitas e a lista dos certificados.

\subsection{Terminologia}\label{sec:intro-terms}

Mantivemos em português a maior parte dos termos técnicos. Na primeira ocorrência, damos entre parênteses o
equivalente em inglês, que é o que se encontra na bibliografia. A \cref{tab:termos} reúne os termos menos usuais e o
sentido em que são usados aqui.

\begin{table}[tbp]
\centering
\small
\begin{tabular}{@{}>{\raggedright\arraybackslash}p{0.22\textwidth}>{\raggedright\arraybackslash}p{0.2\textwidth}
  >{\raggedright\arraybackslash}p{0.5\textwidth}@{}}
\toprule
termo & em inglês & sentido neste artigo \\
\midrule
aritmética de bolas & \emph{ball arithmetic} & cada número real é representado por um centro $m$ e um raio $r$, isto é,
  pelo intervalo $[m-r,m+r]$; cada operação devolve uma bola que contém o resultado exato \\
célula (de uma cobertura) & \emph{tile} & cada uma das regiões de uma cobertura em que se verifica uma desigualdade; as
  células podem se sobrepor \\
deflação & \emph{deflation} & perturbação de posto finito que afasta raízes já conhecidas (aqui, as de gauge) \\
espaço do toro & \emph{torus space} & espaço $\ell^2$ ponderado de coeficientes, igual ao $L^2$ de um toro complexo, no
  qual o fundo age ponto a ponto (\cref{sec:setting-realizations}) \\
faixa A, faixa B & \emph{A band, B band} & regiões do plano do expoente $s$ que contêm as raízes dos setores A e B
  (\cref{sec:setting-sectors}) \\
fundo & \emph{background} & a solução de Reiterer e Trubowitz, em torno da qual se lineariza \\
gauge (calibre) & \emph{gauge} & liberdade de coordenadas; uma transformação de gauge soma à perturbação a derivada de
  Lie do fundo ao longo de um campo vetorial (\cref{def:gauge}) \\
imagem numérica & \emph{numerical range} & o conjunto $W(T)=\{\langle Tx,x\rangle:\ x\in\Dom(T),\ \|x\|=1\}$, onde $\Dom(T)$
  é o domínio de $T$ \\
janela, parte distante & \emph{window, far part} & os níveis intermediário e externo da referência usada na contagem
  (\cref{sec:counting-reference}) \\
localização rigorosa; bola ou intervalo certificado & \emph{enclosure} & um conjunto que comprovadamente contém o valor
  exato \\
recuperação do raio & \emph{radius recovery} & o argumento que leva um autovetor de um espaço de raio menor para o de
  raio maior (\cref{lem:recovery}) \\
setor A, setor B & \emph{sector A, sector B} & os perfis com a simetria do fundo e os com a simetria oposta
  (\cref{sec:setting-sectors}) \\
sistema aumentado & \emph{bordered system} & a equação de autovalor acrescida de uma condição de normalização, com o
  autovalor como incógnita (\cref{sec:roots-nk}) \\
testemunho & \emph{witness} & uma cota inferior que prova que um autovetor viola os vínculos (\cref{prop:witness}) \\
vínculos & \emph{constraints} & as equações do sistema que ficam fora da parte de evolução: $F^\sharp=0$ no sistema de
  Reiterer e Trubowitz e $C(s)h=0$ na linearização (\cref{sec:setting-system}) \\
\bottomrule
\end{tabular}
\caption{Termos técnicos usados neste artigo.}
\label{tab:termos}
\end{table}

\section{O fundo e a formulação de Reiterer--Trubowitz}\label{sec:setting}

Esta seção fixa a notação. As \cref{sec:setting-background,sec:setting-system} recordam as partes
de~\cite{RT2019} que usamos, com a notação de lá sempre que possível. As \cref{sec:setting-linear,sec:setting-sectors,sec:setting-constraints,sec:setting-realizations}
definem os objetos do Teorema Principal: a família linearizada $L(s)$, a aplicação de vínculos $C(s)$, o
operador de propagação dos vínculos $K(s)$, as classes de Floquet e os espaços de funções.
Os números de seção nas referências a~\cite{RT2019} são os da versão 1203.3766v1 do arXiv, cujo fonte usamos
em todo o trabalho.

\subsection{O espaço-tempo de Choptuik}\label{sec:setting-background}

Sejam
\[
  M=\{(u_-,u_+)\in\R^2:\ u_+<0,\ u_+<u_-<-u_+/81\},\qquad
  \xi=\frac{u_+-u_-}{u_++u_-}.
\]
Reiterer e Trubowitz~\cite[Main result]{RT2019} provam a existência de constantes $K,\mu>0$, conhecidas com
$80$ dígitos ($K\approx1.72272620$, $\mu\approx0.16830708$), e de funções analíticas reais
$\phi,\zeta,Q$ em $M$ com $\mu+Q\xi^2>0$, tais que a métrica esfericamente simétrica
\begin{equation}\label{eq:metric}
  \gbar=e^{-2\zeta}\Big(-\frac{2\,(du_-\otimes du_++du_+\otimes du_-)}{\mu^2(u_++u_-)^2}
  +\frac{\xi^2}{(\mu+Q\xi^2)^2}\,g_{S^2}\Big)
\end{equation}
em $M\times S^2$ e o campo escalar $\phibar=\phi$ resolvem
$\operatorname{Ric}_g=2\,d\phi\otimes d\phi$, $\Box_g\phi=0$. A aplicação
$\Theta(u_-,u_+)=e^{-2\pi\mu}(u_-,u_+)$ satisfaz
\[
  \Theta^*\gbar=e^{-2K}\gbar,\qquad \phibar\circ\Theta=-\phibar,
\]
e a curvatura escalar não é identicamente nula; assim o espaço-tempo não é plano e explode no
ponto singular $(u_-,u_+)=(0,0)$, que não pertence a $M$. A fronteira $u_+=u_-$ é o eixo de simetria, uma
singularidade removível das coordenadas. A hipersuperfície nula $u_-=0$ é o cone de luz passado do ponto
singular; $M$ contém uma vizinhança dele.

\emph{Coordenadas.} Usamos as coordenadas $(\tau,\xi)$ de~\cite{RT2019}, definidas por
\begin{equation}\label{eq:tauxi}
  u_\sigma=-(1+\sigma\xi)\,e^{-\mu\tau},\qquad\sigma=\pm.
\end{equation}
Nelas $\Theta$ é a translação $\tau\mapsto\tau+2\pi$, o eixo é $\xi=0$, o cone de luz $u_-=0$ é
$\xi=1$, e a condição $u_-<-u_+/81$ vira $\xi<41/40$. Trocando $\xi\ge0$ pela variável
cartesiana $x\in\R^3$ com $|x|=\xi$, o espaço-tempo passa a ser o aberto
\[
  \Mhat=\R_\tau\times B_{41/40}(0)\subset\R\times\R^3,
\]
no qual $\gbar$ e $\phibar$ são analíticas reais, inclusive no eixo $x=0$ e através do cone de luz
$|x|=1$~\cite{RT2019}. Todas as perturbações abaixo são definidas em $\Mhat$.

\emph{Tempo logarítmico.} Como $\Theta$ reescala os comprimentos por $e^{-K}$ e age por $\tau\mapsto\tau+2\pi$,
o tempo $T=K\tau/2\pi$ é o tempo logarítmico usual do colapso crítico, no qual $\Theta$ é uma
translação por $K$. O campo escalar volta a si mesmo sob $\Theta^2$, de modo que o período de eco é
$\Delta=2K=3.44545240\ldots$, de acordo com o valor numérico $3.445452402(3)$ de
Mart\'in-Garc\'ia e Gundlach~\cite{MartinGarciaGundlach2003}, como já observado em~\cite{RT2019}.

\subsection{Variáveis de referencial e o sistema}\label{sec:setting-system}

Reiterer e Trubowitz escrevem as equações como um sistema de primeira ordem, quadraticamente não linear, para a
constante $\mu$ e quatro funções $\omega=(\omega_1,\omega_2,\omega_3,\omega_4)$ de $(\tau,\xi)$. Com os
campos vetoriais nulos radiais
\[
  D^\sigma=\sigma\,\partial_\tau+\mu(1+\sigma\xi)\,\partial_\xi,\qquad\sigma=\pm,
\]
e a reflexão $(Pf)(\tau,\xi)=f(\tau,-\xi)$, elas são
\begin{equation}\label{eq:omega}
\begin{aligned}
  \omega_1&=2\xi Q+(D^-+D^+)\big(\zeta+\log(\mu+\xi^2Q)\big),&
  \omega_2^\sigma&=D^\sigma\big(\zeta+\log(\mu+\xi^2Q)\big)-\sigma\mu,\\
  \omega_3^\sigma&=D^\sigma\zeta-\sigma\mu,&
  \omega_4^\sigma&=D^\sigma\phi,
\end{aligned}
\end{equation}
com $\omega_i^-=\omega_i$ e $\omega_i^+=-P\omega_i$ para $i=2,3,4$~\cite[\S2]{RT2019}. As equações de
Einstein--campo escalar equivalem à anulação de dez funções explícitas
$\Omega_i^\sigma(\mu,\omega)$, $i\in\{1,2,2\sharp,3,4\}$, cada uma da forma
\[
  \Omega_i^\sigma=\xi\,(D^\sigma+\sigma\mu)(\text{uma componente de }\omega)
  +\mu\,(\text{linear em }\omega)+\xi\,(\text{quadrática em }\omega).
\]
A simetria de reflexão as reduz às cinco equações $\Omega^+=0$, postas no cilindro
$(\R/4\pi\Z)\times(-\xi_*,\xi_*)$, $\xi_*>1$, para funções com as simetrias
\begin{equation}\label{eq:symmetries}
  \omega_1,\omega_2,\omega_3\ \text{$2\pi$-periódicas},\qquad \omega_4\ \text{$2\pi$-antiperiódica},\qquad
  P\omega_1=-\omega_1.
\end{equation}

\emph{Séries.} Cada $\omega_i$ é expandida em série de Fourier em $\tau$ com período $4\pi$ e em série de
Chebyshev em $\xi$:
\begin{equation}\label{eq:series}
  v(\tau,\xi)=\sum_{m,n\in\Z}v_{mn}\,e^{im\tau/2}e^{in\theta},\qquad \xi=\cos\theta,\qquad
  v_{m,-n}=v_{mn}.
\end{equation}
As simetrias~\eqref{eq:symmetries} dizem que $\omega_1,\omega_2,\omega_3$ contêm apenas $m$ par,
$\omega_4$ apenas $m$ ímpar, e $\omega_1$ apenas $n$ ímpar. Nessas coordenadas $\partial_\tau$ age por
$im/2$, a reflexão por $(-1)^n$, a multiplicação por $\xi$ pela média dos vizinhos $n\pm1$, e o
produto de funções pela convolução dos coeficientes.

\emph{O sistema selecionado e os vínculos.} Reiterer e Trubowitz separam as cinco equações
$\Omega^+=0$ numa parte com tantas equações quanto incógnitas e num resto, os vínculos, por meio de dois
operadores de seleção $S$ e $S^\sharp$: $\Omega^+=0$ se e só se $S\Omega^+=0$ e $S^\sharp\Omega^+=0$.
A parte selecionada é
\begin{equation}\label{eq:F}
  F(\mu,\omega):=S\,\Omega^+(\mu,\omega)
  =O_\mu\,\omega+\mu\,S\Gamma_1\omega+S\Xi\,\Gamma_2(\omega,\omega),
\end{equation}
onde:
\begin{itemize}[leftmargin=*]
\item $O_\mu=\partial_\tau+\mu D_O$, com $D_O=\operatorname{diag}(\xi\partial_\xi+1,\
  (1+\xi)\partial_\xi+1,\ (1+\xi)\partial_\xi+1,\ (1+\xi)\partial_\xi+1)$, é a parte livre;
\item $\Gamma_1$ é linear e $\Gamma_2$ é bilinear simétrica; ambas são construídas com reflexões e
  produtos pontuais, sem derivadas;
\item $\Xi$ é a multiplicação por $\xi$, e $S$ contém a divisão regularizada por $\xi$,
  $R_\xi f=(f-f|_{\xi=0})/\xi$.
\end{itemize}
As fórmulas explícitas estão em~\cite[\S3]{RT2019}. Escrevemos $F^\sharp(\mu,\omega):=S^\sharp\Omega^+(\mu,\omega)$
para os vínculos.

\emph{O fundo.} Reiterer e Trubowitz constroem uma solução
$(\mu,\omegabar)$ de $F=0$, $F^\sharp=0$ como ponto fixo de uma contração num espaço $\ell^1$ ponderado
de coeficientes, com pesos $(65/64)^{|m|}(5/4)^{|n|}$, perto de dados numéricos explícitos, e a partir dela o
espaço-tempo da \cref{sec:setting-background}~\cite[\S\S4--5]{RT2019}. A solução é normalizada pela
condição $\operatorname{Re}(\omegabar_4)_{10}=0$, que fixa a invariância por translação em $\tau$. As cotas
finais de lá são
\begin{equation}\label{eq:RTball}
  |\mu-(\mu_A+\mu_B)|\le2^{-277},\qquad \|\omegabar-(\omega_A+\omega_B)\|\le2^{-277},\qquad
  \|\omega_B\|\le2^{-25},
\end{equation}
onde $(\mu_A,\omega_A)$ e $(\mu_B,\omega_B)$ são dados explícitos publicados com~\cite{RT2019}, e a
norma é a norma $\ell^1$ ponderada de lá. A maior parte das nossas cotas usa apenas $(\mu_A,\omega_A)$ e a
consequência $\|\omegabar-\omega_A\|\le2^{-25}+2^{-277}$.

\subsection{A família linearizada}\label{sec:setting-linear}

Uma perturbação do tipo de Floquet é $\omega=\omegabar+\epsilon\,e^{s\tau}h$, onde $h$ é um multicampo
$4\pi$-periódico. Aqui $\mu$ fica fixo, assim como as coordenadas $(\tau,\xi)$: perturbações de $\mu$
moveriam as coordenadas nulas~\eqref{eq:tauxi}. Que isso não restringe a generalidade para os modos físicos é
parte da \cref{sec:physical}. A normalização da fase não é imposta a $h$.

\begin{definition}[A família linearizada]\label{def:L}
Para $s\in\C$ e um multicampo $4\pi$-periódico $h$ com $Ph_1=-h_1$, ponha
\[
  L(s)h:=e^{-s\tau}\,D_\omega F(\mu,\omegabar)\big[e^{s\tau}h\big],\qquad
  C(s)h:=e^{-s\tau}\,D_\omega F^\sharp(\mu,\omegabar)\big[e^{s\tau}h\big].
\]
\end{definition}

\begin{lemma}[$L$ e $C$ são afins em $s$]\label{lem:affine}
$L(s)=L(0)+s\,I$ e $C(s)=C_0+s\,C_1$, onde
\[
  L(0)=J+B,\qquad J=O_\mu,\qquad B=\mu\,S\Gamma_1+2\,S\Xi\,\Gamma_2(\omegabar,\cdot\,),
\]
e $C_0,C_1$ são dados explicitamente no \cref{app:formulas}.
\end{lemma}

\begin{proof}
Em~\eqref{eq:F}, a única derivada em $\tau$ é a de $O_\mu$, com coeficiente um. Os operadores
$\Gamma_1$, $\Gamma_2(\omegabar,\cdot)$, $\Xi$ e $S$ agem em $\xi$ e por produtos pontuais, de modo que
comutam com a multiplicação por $e^{s\tau}$, e $e^{-s\tau}\partial_\tau e^{s\tau}=\partial_\tau+s$. Em
$\Omega^+=\Xi H_\mu\omega+\mu\Gamma_1\omega+\Xi\Gamma_2(\omega,\omega)$ a derivada $\partial_\tau$ aparece
apenas em $H_\mu=H_1\partial_\tau+(\text{termos em }\xi)$, onde $H_1$ é uma matriz constante de identidades e
reflexões~\cite[\S\S2--3]{RT2019}. Portanto o mesmo cálculo dá $C(s)=C_0+sC_1$ com
$C_1=S^\sharp\Xi H_1$.
\end{proof}

Uma raiz de $L$ é um ponto $s_0$ no qual $L(s_0)$ tem núcleo não trivial nas realizações da
\cref{sec:setting-realizations}. Pelo \cref{lem:affine}, as raízes de $L$ são os pontos $s_0$ com
$-s_0\in\sigma(L(0))$, e uma cadeia de Jordan de $L$ em $s_0$, isto é, $L(s_0)h_0=0$ e
$L(s_0)h_j=-h_{j-1}$, é uma cadeia de Jordan do operador $L(0)$ no autovalor $-s_0$. A
\emph{multiplicidade algébrica} $\mult(s_0)$ é a dimensão do autoespaço generalizado
$\bigcup_k\ker(L(0)+s_0)^k$. Ela é finita, e as raízes são isoladas, pela \cref{prop:fredholm}.

\subsection{Classes de Floquet e os dois setores}\label{sec:setting-sectors}

Um multicampo $h$ é $4\pi$-periódico, e portanto contém todos os índices de Fourier $m$; o fundo contém
apenas os índices permitidos por~\eqref{eq:symmetries}. Assim, o espaço dos perfis se decompõe como
$X=X_A\oplus X_B$:
\begin{itemize}[leftmargin=*]
\item \emph{setor A}, com a simetria do fundo: $h_1,h_2,h_3$ com $m$ par e $h_4$ com
  $m$ ímpar;
\item \emph{setor B}, com a simetria oposta: $h_1,h_2,h_3$ com $m$ ímpar e $h_4$ com $m$ par.
\end{itemize}
Como $\omegabar$ pertence ao setor A, os dois setores são invariantes por $L(s)$ e $C(s)$; escrevemos
$L_A,L_B$ e $C_A,C_B$ para as restrições. Em termos do meio período $\tau\mapsto\tau+2\pi$, o setor A
consiste nos perfis com a mesma simetria do fundo, e o setor B nos de simetria oposta.

A multiplicação por $e^{i\tau/2}$ é o deslocamento $m\mapsto m+1$ do índice de Fourier. Ela preserva os
espaços da \cref{sec:setting-realizations} a menos de um fator constante, comuta com toda operação
em $\xi$ e com a multiplicação por funções de $\tau$, e conjuga $\partial_\tau$ em
$\partial_\tau+i/2$. Portanto
\begin{equation}\label{eq:shift}
  e^{-i\tau/2}\,L(s)\,e^{i\tau/2}=L(s+\tfrac{i}{2}),\qquad
  e^{-i\tau/2}\,C(s)\,e^{i\tau/2}=C(s+\tfrac{i}{2}),
\end{equation}
e o deslocamento troca os dois setores.

\begin{lemma}[Redução ao setor A]\label{lem:sectorB}
Para todo $s\in\C$,
\[
  \mult_L(s)=\mult_{L_A}(s)+\mult_{L_A}(s+\tfrac{i}{2}),
\]
e $L_A$ é $i$-periódico: $\mult_{L_A}(s+i)=\mult_{L_A}(s)$. A conjugação~\eqref{eq:shift} leva as
cadeias de Jordan de $L_B$ em $s$ sobre as de $L_A$ em $s+\tfrac i2$, e leva as cadeias que satisfazem os
vínculos em cadeias que satisfazem os vínculos.
\end{lemma}

\begin{proof}
A decomposição $X=X_A\oplus X_B$ reduz $L(s)$, de modo que as multiplicidades se somam. O deslocamento de um
índice de Fourier leva $X_B$ sobre $X_A$ e, por~\eqref{eq:shift}, conjuga $L_B(s)$ a $L_A(s+\tfrac i2)$ e $C_B(s)$
a $C_A(s+\tfrac i2)$; a conjugação preserva núcleos, cadeias de Jordan e seus comprimentos. O deslocamento de dois
índices leva $X_A$ nele mesmo e dá a $i$-periodicidade. Os detalhes para as realizações usadas abaixo,
nas quais o deslocamento é uma isometria a menos do fator $\kappa_1$, estão no \cref{app:lemmas}.
\end{proof}

Consequentemente, as raízes de $L$, contadas por classe de Floquet $s+\tfrac{i}{2}\Z$, correspondem às raízes de
$L_A$ contadas módulo $i\Z$. Todos os cálculos são feitos para $L_A$ na faixa
\[
  \widetilde Q=\{-\tfrac14<\Im s\le\tfrac34\},
\]
cuja parte inferior $\{-\tfrac14<\Im s\le\tfrac14\}$ contém as raízes do setor A (a \emph{faixa A}, \emph{A band}) e cuja metade
superior contém as raízes do setor B, deslocadas por $\tfrac i2$ (a \emph{faixa B}, \emph{B band}). Por exemplo, a raiz $\sB$
do Teorema Principal aparece como $\sB+\tfrac i2$ na faixa B.

\emph{Realidade.} Os coeficientes de $\omegabar$ satisfazem a condição de realidade
$v_{-m,-n}=\overline{v_{mn}}$~\cite{RT2019}; portanto $\overline{L_A(\bar s)\,\bar h}=L_A(s)h$ para a
conjugação correspondente dos coeficientes, e as raízes de $L_A$ são simétricas por $s\mapsto\bar s$.
Para $L_B$ o mesmo argumento, combinado com o deslocamento, dá a simetria $s\mapsto\bar s+i$ na faixa B
(\cref{app:lemmas}). São essas simetrias que tornam reais as quatro raízes do Teorema Principal.

\subsection{Propagação dos vínculos}\label{sec:setting-constraints}

As equações $\Omega^+=0$ não são independentes. Reiterer e Trubowitz mostram que, se $F(\mu,\omega)=0$,
então $\omega^\sharp=F^\sharp(\mu,\omega)$ satisfaz o sistema linear homogêneo
\begin{equation}\label{eq:sharp}
  O_\mu^\sharp\,\omega^\sharp+\mu\,R_\xi\Gamma_1^\sharp\,\omega^\sharp+\Gamma_2^\sharp(\omega,\omega^\sharp)=0,
\end{equation}
onde $O_\mu^\sharp=\partial_\tau+\mu\operatorname{diag}(\xi\partial_\xi+1,(1+\xi)\partial_\xi+1)$, e
$\Gamma^\sharp_1,\Gamma^\sharp_2$ são de novo algébricos~\cite[\S3]{RT2019}. Linearizar essa identidade em
$(\mu,\omegabar)$, sem supor $F(\mu,\omega)=0$, dá o lema seguinte, provado no
\cref{app:lemmas}.

\begin{lemma}[Propagação dos vínculos]\label{lem:KCML}
Sejam
\[
  K(s)=J^\sharp+B^\sharp+s\,I,\qquad J^\sharp=O^\sharp_\mu,\qquad
  B^\sharp=\mu\,R_\xi\Gamma_1^\sharp+\Gamma_2^\sharp(\omegabar,\cdot\,).
\]
Existe uma família afim explícita $M(s)=M_0+sM_1$ (\cref{app:formulas}) tal que
\begin{equation}\label{eq:KCML}
  K(s)\,C(s)=M(s)\,L(s)
\end{equation}
como operadores fechados do domínio de $L$ em $Y(a,b)$ no domínio de $K$ em $Y(a',b')$, para todos os pesos
$1<a'<a$, $1<b'<b$ (\cref{sec:setting-realizations}).
\end{lemma}

Em particular, se $L(s)h=0$ e $K(s)$ é injetivo no seu domínio no espaço menor, então $C(s)h=0$:
um autovetor de $L$ satisfaz os vínculos em todo $s$ que não é raiz de $K$. A perda de
raio de $(a,b)$ para $(a',b')$ é inofensiva, porque todos os nossos certificados para $K$ são formulados no
espaço menor $(129/128,9/8)$.

\subsection{Realizações e a propriedade de Fredholm}\label{sec:setting-realizations}

Usamos duas famílias de espaços de coeficientes. Para pesos $a,b>1$:
\begin{itemize}[leftmargin=*]
\item $Y(a,b)$ é o espaço $\ell^1$ ponderado de~\cite{RT2019}, com norma
  $\sum_i\eta_i\sum_{m,n}a^{|m|}b^{|n|}|(h_i)_{mn}|$; os pesos das componentes $\eta_i>0$ dão normas
  equivalentes e só afetam constantes (\cref{app:lemmas}). Escrevemos $\Yp=Y(65/64,5/4)$ para o espaço da
  prova de existência.
\item $\mathbb T(a,b)$, o \emph{espaço do toro} (\emph{torus space}), é o espaço $\ell^2$ ponderado com índice de Fourier
  \emph{com sinal},
  \[
    \|h\|^2_{\mathbb T(a,b)}=\sum_i\sum_{m,n\in\Z}a^{2m}b^{2n}\,|(h_i)_{mn}|^2 .
  \]
  Ele é o $L^2$ do toro $\{|w|=a\}\times\{|z|=b\}$ nas variáveis $w=e^{i\tau/2}$, $z=e^{i\theta}$,
  $\xi=(z+z^{-1})/2$. Nele, a multiplicação por uma função analítica, a reflexão $P$ e a divisão por
  $\xi$ agem ponto a ponto, o que torna várias cotas ótimas (\cref{sec:counting}).
\end{itemize}
A inclusão $Y(a,b)\subset\mathbb T(a,b)$ tem norma no máximo um.

\begin{proposition}[Propriedade de Fredholm]\label{prop:fredholm}
Em cada um dos espaços $Y(a,b)$ e $\mathbb T(a,b)$ com $1\le a\le65/64$ e $1<b\le5/4$:
\begin{enumerate}[label=\textup{(\roman*)},leftmargin=*]
\item $J$ é fechado no seu domínio maximal $\Dom$, que coincide com o fecho dos polinômios trigonométricos
  na norma do gráfico, e $J+z$ é invertível para $z\ge0$ real, com inversa compacta;
\item $B$ é limitado, de modo que $L(s)=J+B+s$ é fechado no domínio fixo $\Dom$ para todo $s$;
\item com $Q=(J+z_0)^{-1}$ para um $z_0>0$ fixo, a família $L(s)Q=I+(B+s-z_0)Q$ é uma família
  analítica de operadores de Fredholm de índice zero, invertível para $s$ real grande.
\end{enumerate}
Consequentemente, as raízes de $L$ são isoladas e têm multiplicidade algébrica finita. O mesmo vale para
$K(s)$ nos espaços $\sharp$ correspondentes.
\end{proposition}

A prova está no \cref{app:lemmas}. O ponto-chave de~(i) é que as soluções homogêneas de $J+z$ em cada
modo de Fourier são múltiplos de $(1+\xi)^{-1-(z+im/2)/\mu}$ ou de $\xi^{-1-(z+im/2)/\mu}$, que são singulares em
$\xi=-1$ ou em $\xi=0$, dentro da elipse de Bernstein; assim $J+z$ é injetivo no domínio maximal.

Os certificados das \cref{sec:counting,sec:roots} são calculados no espaço do toro, enquanto a prova de
existência, e com ela a interpretação física, vive em $\Yp$. O lema seguinte mostra que a
escolha não importa na região em que contamos.

\begin{lemma}[Independência da realização]\label{lem:Lreal}
Para $|s|\le 3.1$, os núcleos e as cadeias de Jordan de $L(s)$ em $\Yp$ e em $\mathbb T(65/64,5/4)$
coincidem. Portanto, no disco $|s|\le3.1$, as raízes de $L$ e suas multiplicidades algébricas são as mesmas
nas duas realizações.
\end{lemma}

\begin{proof}[Esboço]
Escreva $H(s)=I+(B+s)Q$ com $Q=J^{-1}$, o mesmo operador nos dois espaços, e separe os coeficientes numa
caixa finita $G=\{|m|<1024,\ n<4096\}$ e no seu complemento $T=I-G$. Nos dois espaços
$\|T(H(s)-I)T\|<1$ para $|s|\le3.1$: a cota é $0.371$ no toro, a partir de uma cota
$\|B\|\le3.9642$ certificada pelo símbolo de $B$ (\cref{sec:counting}), e $0.2666$ em $\Yp$, a partir de uma cota na
parte do fundo calculada com os dados e a bola publicada de~\cite{RT2019} (\cref{app:recovery}); ambas usam formas fechadas para $\|QT\|$. Se $H(s)x=0$ com $x$ no
espaço do toro, então $Gx$ é um vetor finito, $Tx$ é a única solução no toro de
$THT\,Tx=-THGx$, e a série de Neumann resolve a mesma equação em $\Yp$; pela unicidade, as duas
soluções coincidem, e $x\in\Yp$. As cadeias são tratadas por indução, já que o lado direito
$-h_{j-1}$ já está no espaço forte. A recíproca é a inclusão. Os detalhes, e a verificação
racional das constantes, estão no \cref{app:lemmas}.
\end{proof}

\begin{remark}\label{rem:radius}
O mesmo argumento de \emph{recuperação do raio} (\emph{radius recovery}) mostra mais: um elemento do núcleo, ou de uma cadeia de Jordan, num
espaço de raio menor está automaticamente no espaço de raio maior, para todo $s$
(\cref{sec:physical}, \cref{lem:recovery}). Ele é usado na \cref{sec:physical} para mostrar que os modos físicos,
que só se supõem analíticos numa vizinhança não especificada, produzem autovetores em $\Yp$.
\end{remark}

\section{Resultado principal}\label{sec:main}

Em todo o texto, $(\gbar,\phibar)$ denota a solução das equações de Einstein--campo escalar
$\operatorname{Ric}_{g}=2\,d\phi\otimes d\phi$, $\Box_{g}\phi=0$ construída por Reiterer e
Trubowitz~\cite{RT2019} na variedade esfericamente simétrica $\Mhat$ da \cref{sec:setting}, com
constante de autossimilaridade $K\approx 1.7227$ e constante de reescala nula $\mu\approx 0.16831$ (ambas conhecidas
com $80$ dígitos). Como em~\cite{RT2019}, nós a chamamos de \emph{espaço-tempo de Choptuik}. A autossimilaridade
discreta $\Theta$ age por $\tau\mapsto\tau+2\pi$ nas coordenadas $(\tau,\xi)$ de RT, reescala os comprimentos por
$e^{-K}$ e inverte o sinal de $\phibar$; o período completo do fundo é, portanto, $\Theta^{2}$, isto é,
$\tau\mapsto\tau+4\pi$.

Uma perturbação é \emph{do tipo de Floquet com expoente $s\in\C$} se ela é $e^{s\tau}$ vezes um perfil
$4\pi$-periódico em $\tau$ (\cref{def:physical-mode}). Como $e^{i\tau/2}$ é ele próprio $4\pi$-periódico, o
expoente só está definido módulo $\tfrac{i}{2}\Z$; chamamos as classes de $\C/\tfrac{i}{2}\Z$ de \emph{classes de
Floquet}, e contamos tudo por classe. Em termos do tempo logarítmico $T$, no qual $\Theta$ é a
translação por $K$, o fator $e^{s\tau}$ é $e^{\lambda T}$ com
\begin{equation}\label{eq:lambda}
  \lambda=\frac{2\pi s}{K}.
\end{equation}

A família $L(s)$, $s\in\C$, é a linearização do sistema de Reiterer--Trubowitz em torno do fundo,
agindo sobre perfis de Floquet (\cref{def:L}); $C(s)$ é a aplicação de vínculos linearizada correspondente. Os
modos de Floquet físicos e as transformações de gauge são definidos na \cref{sec:physical}
(\cref{def:physical-mode,def:gauge}).

\begin{mainthm}
Seja $(\gbar,\phibar)$ o espaço-tempo de Choptuik de~\cite{RT2019}. Considere perturbações lineares
esfericamente simétricas do tipo de Floquet, com $L$ realizado em $\Yp$ ou, de modo equivalente na região considerada,
no espaço do toro (\cref{sec:setting-realizations}); modos físicos na classe de regularidade da \cref{def:physical-mode},
suaves em $\tau$ e analíticos reais em $\xi$ até o cone de luz e através dele; e transformações de gauge como na
\cref{def:gauge}.
\begin{enumerate}[label=\textup{(\arabic*)},leftmargin=*]
\item \textbf{Raízes de $L$.} Módulo $\tfrac{i}{2}\Z$, as raízes de $L$ no semiplano $\Re s>0$, contadas
  com multiplicidade algébrica, são exatamente quatro. Cada uma delas é algebricamente simples e tem um
  representante real:
  \begin{center}
  \begin{tabular}{@{}lll@{}}
  \toprule
  raiz & localização certificada & autovetor de $L$ \\
  \midrule
  $s=\mu$ & $\mu=0.168307078963\ldots$ (exato) & $g_*=(1+Z)\omegabar$, gauge puro \\
  $s=\sB$ & $|\sB-0.0414194105|\le 3.7\cdot10^{-5}$ & viola os vínculos \\
  $s=\sK$ & $|\sK-0.401024733|\le 1.72\cdot10^{-4}$ & viola os vínculos \\
  $s=\sphys$ & $\sphys\in[0.73316949,\ 0.73318983]$ & satisfaz os vínculos \\
  \bottomrule
  \end{tabular}
  \end{center}
  No eixo imaginário, a única raiz de $L$ é $s\equiv 0$; ela é algebricamente simples, com autovetor
  $\partial_\tau\omegabar$, a translação infinitesimal da fase.
\item \textbf{Modos físicos.} O espaço dos modos de Floquet físicos com $\Re s>0$, módulo gauge, tem
  dimensão um: ele consiste nos modos com expoente $\sphys$, e não há modos generalizados
  \textup{(}cadeias de Jordan\textup{)} de comprimento $\ge2$. Se o cone de luz passado da singularidade é mantido fixo
  (\cref{def:conefixed}), o quociente pelo grupo menor das transformações de gauge que preservam o cone tem
  dimensão dois: a classe adicional tem expoente $\mu$ e é tangente à família a um parâmetro de
  soluções exatas obtida transladando o ponto singular, todas isométricas a
  $(\gbar,\phibar)$. Ela não é, portanto, uma segunda instabilidade da solução crítica, mas o efeito de
  fixar o instante da singularidade.
\end{enumerate}
\end{mainthm}

Aqui $Z=\mu^{-1}\partial_\tau+\xi\partial_\xi$. O autovetor $g_*$ é um perfil; o modo de Floquet correspondente
$e^{\mu\tau}g_*$ é a derivada de Lie do fundo ao longo de $X_0=\partial_{u_-}+\partial_{u_+}=e^{\mu\tau}(\mu^{-1}\partial_\tau+\xi\partial_\xi)$,
o gerador das translações do ponto singular $(u_-,u_+)=(0,0)$ (\cref{lem:cone}).

\begin{remark}[O expoente crítico]\label{rem:gamma}
Por~\eqref{eq:lambda} e pelo valor de $K$ de~\cite{RT2019}, o único modo físico instável cresce como
$e^{\lambda T}$ com
\[
  \lambda=\frac{2\pi\sphys}{K}\in[2.674040,\ 2.674115],\qquad
  \frac{1}{\lambda}\in[0.373955,\ 0.373966].
\]
O argumento de escala do colapso crítico relaciona um único modo instável ao expoente de escala da massa por
$\gamma=1/\lambda$~\cite{KoikeHaraAdachi1995,Gundlach1997}. Choptuik mediu $\gamma\approx0.37$ nas suas
simulações~\cite{Choptuik1993}. Gundlach calculou o modo instável por teoria de perturbações numérica e encontrou
$\lambda_1=-2.674\pm0.009$, logo $\gamma=-1/\lambda_1=0.374\pm0.001$~\cite{Gundlach1997}; o sinal é o da sua
variável de tempo, na qual o modo que cresce tem expoente negativo, de modo que o seu $-\lambda_1$ corresponde ao nosso
$\lambda$. O nosso intervalo está dentro da barra de erro dele. Essa concordância é uma conferência externa, não parte do
teorema: o argumento de escala diz respeito à evolução não linear, que não estudamos.
\end{remark}

\begin{remark}[O que se afirma e o que não se afirma]\label{rem:scope}
\begin{enumerate}[label=(\alph*),leftmargin=*]
\item O teorema é \emph{espectral}: ele conta modos de Floquet das equações linearizadas. Não afirma
  nada sobre o semigrupo da evolução linearizada, nem sobre a dinâmica não linear.
\item Só são consideradas perturbações esfericamente simétricas. As perturbações não esféricas são discutidas
  na \cref{sec:discussion}.
\item Os modos físicos são tomados na classe da \cref{def:physical-mode}: suaves em $\tau$ e analíticos
  reais em $\xi$, com uma vizinhança complexa uniforme em $\tau$. A dependência apenas suave em $\xi$
  através do cone de luz e no eixo não é coberta (\cref{sec:discussion}).
\item Reiterer e Trubowitz provam a existência de $(\gbar,\phibar)$ perto dos seus dados numéricos; a unicidade
  local não é enunciada formalmente em~\cite{RT2019}. O teorema trata da solução deles.
\item As raízes de $L$ que violam os vínculos não são físicas: as soluções correspondentes da
  parte de evolução do sistema não resolvem as equações de Einstein. A raiz $\mu$ é um efeito de coordenadas
  do cone de luz móvel: o seu modo é gauge puro, e ele sobrevive apenas no quociente pelo grupo menor
  das transformações de gauge que preservam o cone, que é a segunda contagem em~(2).
\end{enumerate}
\end{remark}

\begin{remark}[Modos neutros]\label{rem:neutral}
A parte~(1) no eixo imaginário diz respeito apenas às raízes de $L$. Em $s=0$ há mais duas simetrias exatas
das equações, a mudança de escala global $\delta g=2\gbar$ e o deslocamento $\delta\phi=\mathrm{const}$,
que são perturbações físicas com $\delta\omega=0$ e, portanto, invisíveis para $L$. Os lemas de gauge e de
reconstrução da \cref{sec:physical} são provados apenas para $\Re s>0$, de modo que a parte~(2) não tem análogo no
eixo.
\end{remark}

\begin{remark}[Dependência do cálculo]\label{rem:computation}
A prova é assistida por computador em dois lugares: a existência do fundo com cotas explícitas, que
tomamos de~\cite{RT2019}, e as cotas das \cref{sec:counting,sec:roots}, que são novas. Todas as
últimas são produzidas em aritmética intervalar (Arb~\cite{Arb}) ou em aritmética racional exata, ou como cotas
a posteriori sobre cálculos em ponto flutuante com controle rigoroso do arredondamento. Os certificados, o programa
de verificação e os dados estão disponíveis publicamente (\cref{sec:computation}).
\end{remark}

\subsection{Estrutura da prova}\label{sec:main-structure}
A prova tem três camadas.
\begin{enumerate}[label=(\roman*),leftmargin=*]
\item \emph{Geometria} (\cref{sec:physical}). Todo modo de Floquet físico com $\Re s>0$ pode ser posto, por uma
  transformação de gauge, na forma usada por Reiterer e Trubowitz, e então corresponde a um elemento de
  $\ker L(s)\cap\ker C(s)$; reciprocamente, todo elemento desses vem de um modo físico. Os únicos modos de gauge
  com $\Re s>0$ são os múltiplos de $e^{\mu\tau}g_*$ (\cref{lem:gauge}). Isso reduz a
  parte~(2) à parte~(1), junto com o comportamento dos vínculos em cada espaço de raízes.
\item \emph{Contagem} (\cref{sec:counting}). O número de raízes de $L$ no semiplano direito é
  calculado por uma versão para operadores do teorema de Rouch\'e com uma referência de dimensão finita, após uma
  deflação de posto dois (de Brauer) das duas raízes de gauge $0$ e $\mu$.
\item \emph{Identificação} (\cref{sec:roots}). Cada raiz é localizada, e mostra-se que é simples, por um
  argumento de Newton--Kantorovich, e mostra-se que o seu autovetor ou viola os vínculos (por uma
  cota inferior explícita), ou é gauge puro, ou, no caso de $\sphys$, os satisfaz em todo o espaço de raízes
  (pela injetividade do operador de propagação dos vínculos $K$).
\end{enumerate}

\section{Dos modos físicos ao operador}\label{sec:physical}

A parte~(2) do Teorema Principal é um enunciado sobre soluções das equações de Einstein--campo escalar
linearizadas, enquanto a parte~(1) trata da família $L(s)$. Esta seção prova que, para $\Re s>0$, as duas são
equivalentes (\cref{thm:correspondence}), e reduz a multiplicidade física de uma raiz de $L$ a álgebra linear
de dimensão finita no seu espaço de raízes (\cref{prop:physmult}). As únicas entradas assistidas por computador são
o fato F1 abaixo e a simplicidade da raiz $\mu$, provada na \cref{sec:roots}.

\subsection{Modos físicos e gauge}\label{sec:physical-defs}

Trabalhamos em $\Mhat=\R_\tau\times B_{41/40}(0)$ com as coordenadas cartesianas $(\tau,x)$, $|x|=\xi$, e nas
coordenadas duplamente nulas $u_\pm$ de~\eqref{eq:tauxi} fora do eixo. Um 2-tensor simétrico, esfericamente
simétrico e suave em $\Mhat$ tem a forma
\begin{equation}\label{eq:invariant-tensor}
  \alpha\,d\tau^2+2\beta\,d\tau\,(x\cdot dx)+\gamma\,(x\cdot dx)^2+d\,|dx|^2,
\end{equation}
com coeficientes que são funções de $(\tau,\xi^2)$.

\emph{A forma RT.} O fundo~\eqref{eq:metric} tem componentes nulas $\gbar_{++}=\gbar_{--}$ nas
coordenadas nulas. Uma perturbação $(\delta g,\delta\phi)$ está \emph{na forma RT} se
$\delta g_{++}=\delta g_{--}=0$; então a sua parte bidimensional é conforme a $\gbar_{2D}$, e a
linearização de~\eqref{eq:metric} e~\eqref{eq:omega} com $\mu$ e coordenadas fixos associa a ela
\emph{variáveis RT} $\delta\omega=D\Omega(\gbar,\phibar)[\delta g,\delta\phi]$, onde
$\Omega:(g,\phi)\mapsto\omega$ é a aplicação~\eqref{eq:omega}. Explicitamente, $\delta g=(\delta a/a)\gbar_{2D}+\delta
b\,g_{S^2}$ dá $\delta\zeta=-\delta a/2a$, $\delta W=-W(\delta\zeta+\delta b/2b)$ com $W=\mu+\xi^2Q$,
$\delta Q=\delta W/\xi^2$, e $\delta\omega$ se obtém derivando~\eqref{eq:omega}
(\cref{app:lemmas}).

\begin{definition}[Modo de Floquet físico]\label{def:physical-mode}
Seja $s\in\C$. Um \emph{modo de Floquet físico com expoente $s$} é um par esfericamente simétrico
$\delta=(\delta g,\delta\phi)$ em $\Mhat$ tal que:
\begin{enumerate}[label=\textup{(\roman*)},leftmargin=*]
\item $\delta$ resolve as equações de Einstein--campo escalar linearizadas em $(\gbar,\phibar)$;
\item $(\Theta^2)^*\delta g=e^{4\pi s}e^{-4K}\delta g$ e $(\Theta^2)^*\delta\phi=e^{4\pi s}\delta\phi$;
\item \emph{(regularidade)} os coeficientes $\alpha,\beta,\gamma,d$ de $\delta g$ em~\eqref{eq:invariant-tensor}
  e $\delta\phi$, multiplicados por $e^{-s\tau}$ e pelo fator de fundo $e^{2\zeta}$ onde ele se aplica
  (de modo que são $4\pi$-periódicos por~(ii)), são suaves em $\Mhat$ e pertencem a $C^\infty(\R/4\pi\Z;\mathcal A(E))$
  para alguma vizinhança complexa $E$ de $[-1,1]$ no plano de $\xi$, onde $\mathcal A(E)$ denota as funções
  holomorfas limitadas em $E$.
\end{enumerate}
Um \emph{modo físico generalizado de ordem $k$} é uma solução de~(i) da forma
$e^{s\tau}\sum_{j<k}\tau^j\delta_j$, com perfis $\delta_j$ como em~(iii).
\end{definition}

A condição~(iii) pede suavidade em $\tau$ e analiticidade real em $\xi$ até o cone de luz $\xi=1$ e através
dele, com uma vizinhança complexa uniforme em $\tau$, mas que pode ser arbitrariamente fina. O fator
$e^{-4K}$ em~(ii) é o do fundo, $(\Theta^2)^*\gbar=e^{-4K}\gbar$; de modo equivalente,
$\delta\zeta$, $\delta W/W$ e $\delta\phi$ são $e^{s\tau}$ vezes funções $4\pi$-periódicas.

\begin{definition}[Gauge]\label{def:gauge}
Uma \emph{transformação de gauge} é $\delta\mapsto\delta+\Lie_X(\gbar,\phibar)$, onde
$\Lie_X(\gbar,\phibar)=(\Lie_X\gbar,X(\phibar))$ e $X$ é um campo vetorial complexo em $\Mhat$ que é
\begin{enumerate}[label=\textup{(\alph*)},leftmargin=*]
\item de classe $C^1$ em $\Mhat$, ou
\item contínuo em $\Mhat$, de classe $C^1$ fora do cone de luz, e tal que $\Lie_X(\gbar,\phibar)$ está na
  forma RT com variáveis RT analíticas.
\end{enumerate}
O parâmetro $\mu$ e as coordenadas $(\tau,x)$ ficam fixos. Um modo físico é \emph{gauge puro}
se ele é igual a $\Lie_X(\gbar,\phibar)$ para um tal $X$.
\end{definition}

\begin{definition}[Cone de luz fixo]\label{def:conefixed}
No \emph{cenário de cone fixo}, exige-se que os modos físicos deixem o cone de luz nulo em primeira ordem,
$\delta g_{++}=0$ em $\{u_-=0\}$, e que as transformações de gauge sejam tangentes a ele,
$X(u_-)=0$ em $\{u_-=0\}$.
\end{definition}

Manter $\mu$ fixo é essencial: mudar $\mu$ muda as direções nulas
$\ker(\sigma\,d\xi-\mu(1+\sigma\xi)\,d\tau)$, e uma perturbação na forma RT para um valor de $\mu$ não está na
forma RT para outro. Fixar $\mu$ e as coordenadas faz parte da identificação do fundo.

Para $X$ tal que $\Lie_X(\gbar,\phibar)$ está na forma RT, escrevemos
$\delta\omega_X:=D\Omega(\gbar,\phibar)[\Lie_X(\gbar,\phibar)]$. O exemplo básico é
\begin{equation}\label{eq:X0}
  X_0=\partial_{u_-}+\partial_{u_+}=e^{\mu\tau}\big(\mu^{-1}\partial_\tau+\xi\partial_\xi\big),
\end{equation}
o gerador da translação conjunta das coordenadas nulas, que move o ponto singular $(0,0)$ e
o cone de luz. Ele é analítico em $\Mhat$, $\Lie_{X_0}(\gbar,\phibar)$ está na forma RT, e
\begin{equation}\label{eq:gaugemode}
  G_*:=\delta\omega_{X_0}=e^{\mu\tau}g_*,\qquad g_*=(1+Z)\,\omegabar,\qquad Z=\mu^{-1}\partial_\tau+\xi\partial_\xi .
\end{equation}
A identidade~\eqref{eq:gaugemode} é um cálculo direto a partir de~\eqref{eq:omega} (\cref{app:lemmas}). Como
$\Lie_{X_0}(\gbar,\phibar)$ resolve as equações linearizadas, o passo (R4) do \cref{lem:R} abaixo mostra que
$e^{\mu\tau}g_*$ é um modo de Floquet de $L$ com expoente $\mu$: $L(\mu)g_*=0$ e $C(\mu)g_*=0$.

\subsection{Modos de gauge}\label{sec:physical-gauge}

O fato seguinte é conferido em aritmética racional exata a partir dos dados de~\cite{RT2019}.

\begin{lemma}[F1]\label{lem:F1}
$\omegabar_4(\cdot,0)\not\equiv0$ e $\omegabar_4(\cdot,-1)\not\equiv0$. Mais precisamente, os coeficientes de
Fourier de índice $m=1$ satisfazem $|c_1(\omegabar_4(\cdot,0))|\ge0.493046-3\cdot10^{-8}$ e
$|c_1(\omegabar_4(\cdot,-1))|\ge0.097752-3\cdot10^{-8}$.
\end{lemma}

\begin{proof}
Com $\omega(\tau,\xi)=\sum v_{mn}e^{im\tau/2+in\theta}$, $\xi=\cos\theta$, o coeficiente é
$c_1=\sum_nv_{1n}i^n$ em $\xi=0$ e $\sum_nv_{1n}(-1)^n$ em $\xi=-1$. Para os dados $\omega_A$ essas somas são
aproximadamente $0.0729160+0.4930463\,i$ e $-0.0794729+0.0977527\,i$, calculadas exatamente. Por~\eqref{eq:RTball},
e como os pesos da norma são $\ge1$, $|c_1(\omegabar)-c_1(\omega_A)|\le2^{-25}+2^{-277}<3\cdot10^{-8}$.
\end{proof}

\begin{lemma}[Completude do gauge]\label{lem:gauge}
Seja $X$ um campo vetorial como na \cref{def:gauge}\textup{(a)} tal que $\Lie_X(\gbar,\phibar)$ é esfericamente
simétrico e está na forma RT,
e suponha que $\delta\omega_X\neq0$ é do tipo de Floquet com expoente $s$, $\Re s>0$, isto é,
$\delta\omega_X\circ\Theta^2=e^{4\pi s}\delta\omega_X$. Então $e^{4\pi(s-\mu)}=1$ e
$\Lie_X(\gbar,\phibar)=c\,\Lie_{X_0}(\gbar,\phibar)$ para algum $c\in\C\setminus\{0\}$; em particular,
$\delta\omega_X=c\,G_*$. Consequentemente, a média de $X$ sobre $SO(3)$ é $cX_0$, e $X-cX_0$ é um campo de Killing de
$\gbar$ que anula $\phibar$, por exemplo uma rotação. O mesmo vale para $X$ como na
\cref{def:gauge}\textup{(b)}.
\end{lemma}

Em palavras: o único modo de gauge com $\Re s>0$ é $e^{\mu\tau}g_*$, com expoente $\mu$, no setor A.

\begin{proof}
\emph{Passo 0 (simetria).} Tomar a média de $X$ sobre $SO(3)$ com respeito à medida de Haar não muda
$\Lie_X(\gbar,\phibar)$, porque o fundo e $\Lie_X(\gbar,\phibar)$ são invariantes. A média $\bar X$ é
$C^1$ e invariante, logo da forma $X^\tau\partial_\tau+X^\xi\partial_\xi$ e tangente ao eixo, e
$\Lie_{\bar X}(\gbar,\phibar)=\Lie_X(\gbar,\phibar)$. Trocamos $X$ por $\bar X$; as conclusões abaixo dizem respeito a
$\bar X$, e $X-\bar X$ tem derivada de Lie do fundo nula. (O campo original não precisa ser
invariante: campos de Killing de rotação podem ser somados sem mudar a perturbação.)

\emph{Passo 1 (transporte).} Em coordenadas nulas, $\gbar_{\pm\pm}=0$ e $\gbar_{+-}\neq0$, de modo que
$(\Lie_X\gbar)_{++}=2\gbar_{+-}\partial_{u_+}X^{u_-}$ e $(\Lie_X\gbar)_{--}=2\gbar_{+-}\partial_{u_-}X^{u_+}$.
A forma RT obriga $\partial_{u_+}X^{u_-}=\partial_{u_-}X^{u_+}=0$ em $M$. Os conjuntos de nível
$\{u_-=v\}\cap M$ e $\{u_+=w\}\cap M$ são intervalos, logo conexos, e assim $X^{u_-}=f_-(u_-)$ e
$X^{u_+}=f_+(u_+)$ em $M$. Como $u_-$ assume todos os valores reais em $M$, inclusive $0$ no cone de luz,
$f_-\in C^1(\R)$; e $f_+\in C^1((-\infty,0))$.

\emph{Passo 2 (eixo).} $X$ é tangente ao eixo $u_+=u_-$, logo $f_+=f_-$ em $(-\infty,0)$. Escreva $f:=f_-$.

\emph{Passo 3 (equivariância).} $\Theta$ preserva a forma RT e o referencial de~\cite{RT2019} tem coeficientes
independentes de $\tau$, de modo que $\Omega(\Theta^*(g,\phi))=\Omega(g,\phi)\circ\Theta$; e $\Omega$ é invariante por
reescalas conformes constantes de $g$. Derivando, e usando $(\Theta^2)^*\gbar=e^{-4K}\gbar$,
$\phibar\circ\Theta^2=\phibar$, obtemos $\delta\omega_X\circ\Theta^2=\delta\omega_{(\Theta^2)^*X}$. Como
$u_\pm\circ\Theta^2=\Lambda u_\pm$ com $\Lambda=e^{-4\pi\mu}$, o campo $(\Theta^2)^*X$ é do mesmo tipo,
com $f$ trocada por $\Lambda^{-1}f(\Lambda\,\cdot)$.

\emph{Passo 4 (injetividade).} Suponha $\delta\omega_X=0$. Então $D^\sigma(X(\phibar))=\delta\omega_4^\sigma=0$
para $\sigma=\pm$, de modo que $X(\phibar)=c_2$ é constante no conjunto conexo $M$. Como $D^+u_+=D^-u_-=0$ e
$D^+u_-=-D^-u_+=2\mu e^{-\mu\tau}$,
\[
  X(\phibar)=\frac{e^{\mu\tau}}{2\mu}\big(f(u_-)\,\omegabar_4^+-f(u_+)\,\omegabar_4^-\big).
\]
No eixo $u_\pm=-e^{-\mu\tau}$ e $\omegabar_4^\pm(\tau,0)=\mp\omegabar_4(\tau,0)$, o que dá
$-\mu^{-1}e^{\mu\tau}f(-e^{-\mu\tau})\,\omegabar_4(\tau,0)=c_2$ para todo $\tau$. Como $\omegabar_4(\cdot,0)$ é
antiperiódica, ela se anula em algum ponto, logo $c_2=0$. Pelo \cref{lem:F1} ela não é identicamente nula e,
sendo analítica real, tem zeros isolados; portanto $f=0$ num subconjunto denso de $(-\infty,0)$ e, por continuidade, em
$(-\infty,0]$. Fora do cone, $\{\xi>1\}\cap M$, temos $u_+<0$, logo $f(u_-)\,\omegabar_4^+=0$ ali. Se
$f(v)\neq0$ para algum $v>0$, então $\omegabar_4^+$ se anula num aberto não vazio, logo identicamente por
analiticidade, e em particular no eixo, contradizendo o \cref{lem:F1}. Assim $f\equiv0$ e $X=0$. (Para
$X$ complexo, aplique isso às partes real e imaginária.)

\emph{Passo 5 (homogeneidade).} Se $\delta\omega_X\circ\Theta^2=e^{4\pi s}\delta\omega_X$, os Passos 3 e 4 dão
$(\Theta^2)^*X=e^{4\pi s}X$, isto é,
\[
  f(\Lambda u)=\rho\,f(u)\quad(u\in\R),\qquad \rho=e^{4\pi(s-\mu)},\quad |\rho|=e^{4\pi(\Re s-\mu)}.
\]

\emph{Passo 6 (regularidade no cone).} $f$ é $C^1$ em $u=0$, porque o cone está dentro de $M$. Iterando,
$f(\Lambda^ju)=\rho^jf(u)$ com $\Lambda^ju\to0$.
\begin{itemize}[leftmargin=*]
\item Se $|\rho|>1$: o lado esquerdo tende a $f(0)$, logo $f(u)=0$.
\item Se $|\rho|=1$ e $\rho\neq1$: $\rho^j$ não converge, logo $f(u)=0$. Se $\rho=1$: $f(u)=f(0)$.
\item Se $\Lambda<|\rho|<1$, isto é, $0<\Re s<\mu$: $f(0)=0$, e
  $f'(0)=\lim_jf(\Lambda^ju)/(\Lambda^ju)=\lim_j(\rho/\Lambda)^jf(u)/u$ com $|\rho/\Lambda|>1$, logo $f(u)=0$.
\end{itemize}
Como $\delta\omega_X\neq0$, $f\not\equiv0$, logo $\rho=1$ e $f$ é constante, isto é, $X=cX_0$.

\emph{Variante (b).} Os Passos 1--5 usam apenas a continuidade de $f$ em $u=0$, e o Passo 6 também, para $\Re s\ge\mu$.
Para $0<\Re s<\mu$: $\delta\phi=X(\phibar)$ é contínua com derivadas analíticas
$D^\sigma\delta\phi=\delta\omega_4^\sigma$, logo analítica. No eixo,
$-\mu^{-1}e^{\mu\tau}f(-e^{-\mu\tau})\omegabar_4(\tau,0)$ é analítica, logo $f$ é analítica em $(-\infty,0)$.
Perto de um ponto $(\tau_0,1)$ do cone com $\omegabar_4^+(\tau_0,1)\neq0$, que existe pelo \cref{lem:F1},
$f(u_-)=\big(2\mu e^{-\mu\tau}\delta\phi+f(u_+)\omegabar_4^-\big)/\omegabar_4^+$ é analítica, logo $f$ é analítica
em $0$. Então $f(\Lambda u)=\rho f(u)$ obriga $a_k(\Lambda^k-\rho)=0$ para os coeficientes de Taylor, de modo que
$f=a_ku^k$ com $s\equiv\mu(1-k)\pmod{\tfrac i2\Z}$, e $\Re s>0$ obriga $k=0$.
\end{proof}

\begin{remark}[O papel do cone de luz]\label{rem:cone-role}
A prova usa que o cone de luz está dentro do domínio e que $X$ é contínuo através dele. Se de $X$ se exigisse
apenas regularidade em $\{\xi<1\}$, então $f(u)=(-u)^{1-s/\mu}$ produziria um ``modo de gauge'' para
todo $s$. Exigir um gauge que preserve o cone significa $f(0)=0$; isso exclui $f$ constante, e então não há
modo de gauge algum com $\Re s>0$.
\end{remark}

O modo de gauge $G_*$ tem um significado geométrico direto.

\begin{lemma}[O cone]\label{lem:cone}
Seja $\psi_\epsilon(u_-,u_+)=(u_-+\epsilon,u_++\epsilon)$ o fluxo de $X_0$, estendido trivialmente às
esferas. Para $\epsilon$ pequeno, $(\gbar_\epsilon,\phibar_\epsilon):=\psi_\epsilon^*(\gbar,\phibar)$ é uma solução
exata das equações de Einstein--campo escalar em subconjuntos compactos de $\Mhat$, isométrica a
$(\gbar,\phibar)$, discretamente autossimilar em torno do ponto singular deslocado $(-\epsilon,-\epsilon)$, e na forma
RT com o mesmo $\mu$ e as mesmas coordenadas. As suas variáveis RT resolvem as dez equações $\Omega_i^\sigma=0$, e
\[
  \Omega(\gbar_\epsilon,\phibar_\epsilon)=\omegabar+\epsilon\,e^{\mu\tau}g_*+O(\epsilon^2)
\]
uniformemente em subconjuntos compactos.
\end{lemma}

\begin{proof}
As equações são invariantes por difeomorfismos. A aplicação $\psi_\epsilon$ preserva $du_\pm$ e o eixo
$\{u_+=u_-\}$, de modo que a parte bidimensional de $\gbar_\epsilon$ continua conforme a $du_-du_+$, e
$\gbar_\epsilon$ tem a forma~\eqref{eq:metric} com $e^{-2\zeta_\epsilon}=e^{-2\zeta\circ\psi_\epsilon}
(u_++u_-)^2/(u_++u_-+2\epsilon)^2$ e $Q_\epsilon$ determinado pelo raio de área; a regularidade no eixo
decorre da regularidade da métrica. A derivada em $\epsilon=0$ é
$D\Omega[\Lie_{X_0}(\gbar,\phibar)]=G_*$, por~\eqref{eq:gaugemode}.
\end{proof}

Assim, $e^{\mu\tau}g_*$ é tangente a uma curva de soluções exatas, todas isométricas ao fundo: a mesma
solução crítica com o instante $T_*$ da singularidade transladado. Com as transformações de gauge da
\cref{def:gauge} ele é gauge puro. No cenário de cone fixo da \cref{def:conefixed} não é, pois $X_0(u_-)=1$, e ele
conta como modo físico; esta é a segunda contagem da parte~(2) do Teorema Principal.

\subsection{Reconstrução: dos modos físicos a $\ker L\cap\ker C$}\label{sec:physical-forward}

Precisamos de um enunciado de recuperação do raio para todo $s$, e não só para $|s|\le3.1$ como no \cref{lem:Lreal}.

\begin{lemma}[Recuperação do raio]\label{lem:recovery}
Sejam $1\le\kappa_1'\le65/64$, $1<\kappa_2'\le9/8$ e $S>0$. Para $|s|\le S$, todo elemento do núcleo de
$L(s)$, e toda cadeia de Jordan de $L$ em $s$, no domínio de $L$ em $Y(\kappa_1',\kappa_2')$ está no
domínio de $L$ em $\Yp$.
\end{lemma}

\begin{proof}[Esboço]
Como no \cref{lem:Lreal}: com $Q=J^{-1}$, $H(s)=I+(B+s)Q$ e uma caixa finita $G$, basta que
$\|T(H(s)-I)T\|<1$ nos dois espaços. No espaço fraco isso vale porque $\|QT\|$ tende a zero com $G$,
uniformemente no peso de Fourier (já que $Q$ preserva $m$), e $\|B\|$ é controlada pela cota
$\beta_+$ do \cref{app:recovery}, calculada com os dados e a bola publicada de~\cite{RT2019}, reponderada de $5/4$
para $\kappa_2'$; no espaço forte vale o mesmo com o próprio $\beta_+$. O peso de Fourier $\kappa_1'=1$ é permitido, porque a convolução pelo fundo tem
norma no máximo $\sum_k|b_k|\le\sum_k\kappa_1^{|k|}|b_k|$. Então a equação do núcleo é resolvida no espaço forte
por uma série de Neumann em $T$, e a unicidade no espaço fraco identifica as duas soluções; as cadeias seguem por
indução. Caixas explícitas, por exemplo $G=2^{17}\times2^{20}$ para $S=200$ e $\kappa_2'=1+2^{-6}$, e a
verificação racional das constantes estão no \cref{app:lemmas}.
\end{proof}

\begin{lemma}[Reconstrução, sentido direto]\label{lem:R}
Seja $\delta$ um modo de Floquet físico com expoente $s$, $\Re s>0$. Então:
\begin{enumerate}[label=\textup{(\arabic*)},leftmargin=*]
\item existe uma transformação de gauge, por um campo vetorial $X$ da mesma regularidade que $\delta$ (em
  particular $C^1$ em $\Mhat$) e do tipo de Floquet, tal que $\delta'=\delta+\Lie_X(\gbar,\phibar)$ está na forma
  RT;
\item as variáveis RT de $\delta'$ são $\delta\omega'=e^{s\tau}h$ com $h$ no domínio de $L$ em $\Yp$,
  $L(s)h=0$ e $C(s)h=0$;
\item $h$ é determinado pela classe de gauge de $\delta$ a menos da direção de gauge
  $\mathcal G_s$, onde $\mathcal G_s=\C\,e^{(\mu-s)\tau}g_*$ se $s\equiv\mu\pmod{\tfrac i2\Z}$ e
  $\mathcal G_s=0$ caso contrário; e $h\in\mathcal G_s$ se e só se $\delta$ é gauge puro;
\item um modo físico generalizado de ordem $k$ dá, do mesmo modo, uma cadeia de Jordan de $L$ em $s$ de comprimento
  $k$ módulo $\mathcal G_s$, que satisfaz os vínculos de jato~\eqref{eq:jet} abaixo.
\end{enumerate}
\end{lemma}

\begin{proof}
\emph{(R1) Fixação do gauge fora da ressonância.} Procuramos $X$ com $(\delta g+\Lie_X\gbar)_{\pm\pm}=0$, isto
é, $\partial_{u_+}X^{u_-}=-\delta g_{++}/(2\gbar_{+-})$ e $\partial_{u_-}X^{u_+}=-\delta g_{--}/(2\gbar_{+-})$.
Com $X^{u_\mp}=e^{(s-\mu)\tau}y_\mp$ e $f_\mp:=-\mu e^{-s\tau}\delta g_{\pm\pm}/\gbar_{+-}$, que são
$4\pi$-periódicas pela \cref{def:physical-mode}(ii) (tanto $\delta g_{\pm\pm}$ quanto $\gbar_{+-}$ escalam por
$e^{8\pi\mu}e^{-4K}$ sob $\Theta^2$, a primeira com o fator extra $e^{4\pi s}$), as equações viram
\[
  T_{\pm1}(s)\,y_\mp=f_\mp,\qquad T_c(s)=\partial_\tau+\mu(\xi-c)\partial_\xi+s-\mu .
\]
Para $\Re s>0$ e $s\notin\mu-\tfrac i2\Z$, a única solução periódica, analítica em $\xi$ perto de $[-1,1]$,
é
\begin{equation}\label{eq:transport}
  y(\tau,\xi)=(\partial_\tau+s-\mu)^{-1}f(\cdot,c)\,(\tau)
  +\int_0^\infty e^{-(s-\mu)t}\,\big(f-f(\cdot,c)\big)\big(\tau-t,\ c+e^{-\mu t}(\xi-c)\big)\,dt,
\end{equation}
onde $(\partial_\tau+s-\mu)^{-1}$ é o multiplicador de Fourier $1/(s-\mu+im/2)$. A integral converge também
para $0<\Re s<\mu$, porque o integrando se anula em $\xi=c$ e é, portanto, $O(e^{-\Re s\,t})$; o
caminho característico fica numa vizinhança convexa de $[-1,1]$. A unicidade segue expandindo uma
solução homogênea em potências de $\xi-c$: o coeficiente de grau $n$ resolve
$(\partial_\tau+s+\mu(n-1))y_n=0$, o que obriga $y_n=0$. A fórmula comuta com $\partial_\tau$, de modo que $y$
tem a regularidade de $f$. No eixo, a regularidade cartesiana de $\delta g$ faz de $f_-$ e $Pf_+$ a
mesma função analítica, e a unicidade dá $y_+=Py_-$, o que significa que $X$ é regular no eixo.

\emph{(R2) Ressonância.} Seja $s_0=\mu-im_0/2$ com $\Re s_0=\mu>0$. Em $\xi=1$ a equação para $y_-$ fica
$(\partial_\tau-im_0/2)\,y_-(\cdot,1)=f_-(\cdot,1)$, de modo que existe solução periódica se e só se o
coeficiente $\kappa=(f_-(\cdot,1))_{m_0}$ se anula; nesse caso, a solução é única a menos da soma de
$e^{im_0\tau/2}$, isto é, a menos da soma de um múltiplo de $X_0$. Suponha $\kappa\neq0$. Então
$y_-=y_{\mathrm{reg}}+\kappa\tau e^{im_0\tau/2}$, e a mesma constante aparece em $y_+$ porque
$f_+(\cdot,-1)=f_-(\cdot,1)$; a regularidade no eixo exige a mesma constante nas duas componentes. O
campo de gauge é $X=X_{\mathrm{reg}}+\kappa\tau X_0$, e
\[
  \delta'=A+\tau B,\qquad B=\kappa\Lie_{X_0}(\gbar,\phibar),\qquad
  A=\delta+\Lie_{X_{\mathrm{reg}}}(\gbar,\phibar)+2\kappa\,d\tau\odot X_0^\flat,
\]
ambos na forma RT e do tipo de Floquet. Como $\delta\omega$ é de primeira ordem com $D^\sigma\tau=\sigma$, as variáveis
RT são $\delta\omega'=e^{s_0\tau}(h_1+\tau h_0)$ com $h_0=\kappa e^{im_0\tau/2}g_*\neq0$, e pelo
\cref{lem:affine} as equações linearizadas dão $L(s_0)h_0=0$ e $L(s_0)h_1=-h_0$. Multiplicando por
$e^{-im_0\tau/2}$ (\cref{eq:shift}) e projetando no setor A, onde está $g_*$, obtém-se uma cadeia de Jordan
de $L_A$ de comprimento dois em $s=\mu$, no domínio em $\Yp$ pelo \cref{lem:recovery}. Isso contradiz a
simplicidade algébrica da raiz $\mu$ (\cref{prop:mu}). Portanto $\kappa=0$, e a fixação do gauge existe e é
única a menos de múltiplos de $X_0$.

\emph{(R3) Variáveis RT.} Para $\delta'$ na forma RT, os coeficientes~\eqref{eq:invariant-tensor} dão
$\delta b=d\,\xi^2$, $\delta\zeta=-\tfrac12\mu^2e^{2\zeta}(\gamma\xi^2+d)$ e
$\delta Q=-\tfrac12We^{2\zeta}\big(d(2\mu Q+\xi^2Q^2)-\mu^2\gamma\big)$, que são pares e analíticos em $\xi$.
Portanto $\delta\omega'=e^{s\tau}h$ com $h$ de classe $C^\infty(\R/4\pi\Z;\mathcal A(E'))$ para uma vizinhança
$E'$ um pouco menor; a codificação $\delta\omega_i^+=-P\delta\omega_i^-$ é automática.

\emph{(R4) Equações.} As definições~\eqref{eq:omega} implicam, após a linearização, as identidades
$\Omega_1^\sigma=\Omega_2^\sigma+\Omega_{2\sharp}^\sigma$ e $\Omega_j^-=\Omega_j^+$ para $j=2,3,4$, e as
fórmulas de curvatura de~\cite[\S2]{RT2019} exprimem as equações de Einstein--campo escalar linearizadas pelas
combinações restantes dos $\Omega$. Como $\delta'$ resolve as equações linearizadas, todos os $\delta\Omega_i^\sigma$
se anulam fora do eixo e, por continuidade, no eixo. Em particular, $L(s)h=0$ e $C(s)h=0$.

\emph{(R5) O espaço espectral.} Os coeficientes de Fourier de $h$ decaem mais rápido que qualquer potência de $|m|$, e
os seus coeficientes de Chebyshev decaem como $\varrho^{-n}$, para o parâmetro $\varrho>1$ de uma elipse de Bernstein
contida em $E'$. Portanto $h\in Y(1,\kappa_2')$ para $1<\kappa_2'<\min(\varrho,9/8)$, e $Jh=-(B+s)h$ está no mesmo
espaço, de modo que $h$ está no domínio de $L$ ali. Pelo \cref{lem:recovery}, $h$ está no domínio em $\Yp$; em particular, $e^{s\tau}h$ se estende analiticamente a
$|\xi|<41/40$. Além do cone de luz, para $1<\xi<41/40$, as duas famílias de características $D^\pm$ satisfazem
$d\xi/d\tau>0$, de modo que as hipersuperfícies $\{\xi=\mathrm{const}\}$ são do tipo espaço, e uma solução do sistema
linearizado ali é determinada pelos seus dados de Cauchy em $\{\xi=1+\epsilon\}$, uniformemente no círculo de $\tau$ do
perfil periódico. Como $\delta\omega'$ e $e^{s\tau}h$ coincidem perto de $[-1,1]$ e ambas são soluções suaves, elas
coincidem em todo $\Mhat$. O mesmo se aplica ao próprio $\delta'$ em (R6).

\emph{(R6) Unicidade módulo gauge.} Se $\delta\omega'=0$, então~\eqref{eq:omega} dão $\delta Q=0$ e
$\delta\zeta,\delta\phi$ constantes; uma constante não é do tipo de Floquet com $\Re s>0$, logo $\delta'=0$ e $\delta$
é gauge puro. Se $\delta\omega'=c\,G_*$ (no rótulo $s$), então $\delta+\Lie_{X-cX_0}(\gbar,\phibar)$ tem
$\delta\omega=0$ e vale a mesma conclusão. Duas fixações de gauge diferem por um campo admissível, cujo
$\delta\omega$ está em $\mathcal G_s$ pelo \cref{lem:gauge}.

\emph{(R7) Cadeias.} Para um modo generalizado, o transporte é resolvido ordem a ordem,
$T_c(s)y_j=f_j-y_{j-1}$, com campos de gauge polinomiais em $\tau$. Fora da ressonância toda ordem tem solução.
Na ressonância, a condição de compatibilidade em cada ordem é forçada como em (R2): um termo secular a mais
alongaria uma cadeia de Jordan de $L$ em $\mu$, que é simples. Os passos (R3)--(R5) se aplicam a cada perfil e dão
\[
  \delta\omega'=e^{s\tau}\sum_{j=0}^{k-1}\frac{\tau^j}{j!}\,h_{k-1-j}.
\]
Como $\partial_\tau$ entra nas equações selecionadas e nos vínculos apenas pelos coeficientes constantes
$1$ e $C_1$ (\cref{lem:affine}), o operador selecionado linearizado leva $e^{s\tau}\tau^jh$ em
$e^{s\tau}\big(\tau^jL(s)h+j\tau^{j-1}h\big)$, e a aplicação de vínculos linearizada o leva em
$e^{s\tau}\big(\tau^jC(s)h+j\tau^{j-1}C_1h\big)$. Agrupando as potências de $\tau$, as equações linearizadas viram
$L(s)h_i=-h_{i-1}$ e os vínculos $C(s)h_i+C_1h_{i-1}=0$ ($h_{-1}=0$): $(h_0,\dots,h_{k-1})$ é uma cadeia de Jordan de
$L$ em $s$, no domínio em $\Yp$, que satisfaz os vínculos de jato~\eqref{eq:jet}. Para $k=2$ esta é a forma
$e^{s\tau}(h_1+\tau h_0)$ de (R2). Se
$h_0\in\mathcal G_s$, então, por (R6), a ordem mais baixa é gauge puro, e a cadeia encurta módulo gauge.
\end{proof}

\subsection{A recíproca}\label{sec:physical-converse}

\begin{lemma}[Reconstrução, sentido recíproco]\label{lem:V}
Seja $\Re s>0$, e seja $h$ um multicampo $4\pi$-periódico com $Ph_1=-h_1$, no domínio de $L$ em
$\Yp$, com $L(s)h=0$ e $C(s)h=0$. Então existe um único $\delta=(\delta g,\delta\phi)$ na forma RT cujas variáveis
RT são $e^{s\tau}h$, e $\delta$ é um modo de Floquet físico com expoente $s$, analítico real em
$\Mhat$, inclusive no eixo e no cone de luz. A aplicação $h\mapsto\delta$ é linear e injetiva.
\end{lemma}

\begin{proof}
\emph{As dez equações.} As componentes selecionadas de $\delta\Omega^+$ se anulam em $\xi=0$ para todo $h$ com $h_1$
ímpar, porque as suas partes principal e quadrática carregam um fator $\xi$; assim, nada se perde na divisão
$R_\xi$ contida em $S$. Portanto $L(s)h=0$ dá $\tfrac12(1+P)\delta\Omega_1^+=\delta\Omega_2^+=\delta\Omega_3^+=
\delta\Omega_4^+=0$, e $C(s)h=0$ dá $\tfrac12(1-P)\delta\Omega_1^+=\delta\Omega_{2\sharp}^+=0$. Assim
$\delta\Omega^+=0$, e $P\delta\Omega_i^\sigma=-\delta\Omega_i^{-\sigma}$ dá $\delta\Omega^-=0$.

\emph{Potenciais.} Procuramos $(\delta Q,\delta\zeta,\delta\phi)=e^{s\tau}(q,z,p)$ com
$\delta\omega=e^{s\tau}h$. Por~\eqref{eq:omega} e~\eqref{eq:metric}, isso significa
$D^\sigma_s z=h_3^\sigma$, $D_s^\sigma p=h_4^\sigma$ e $D_s^\sigma(z+\delta\log W)=h_2^\sigma$, com
$D_s^\sigma=D^\sigma+\sigma s$, junto com a relação algébrica para $q$ vinda de $\omega_1$. Com
$a_i=\tfrac12[(1-\xi)h_i^+-(1+\xi)h_i^-]$ e $b_i=(h_i^++h_i^-)/(2\mu)$, a equação $D_s^\sigma z=h_i^\sigma$ para
os dois sinais equivale a $(\partial_\tau+s)z=a_i$ e $\partial_\xi z=b_i$. Um cálculo direto dá
\[
  -2\mu\xi\big((\partial_\tau+s)b_i-\partial_\xi a_i\big)=\delta\Omega_j^--\delta\Omega_j^+,\qquad (i,j)=(3,3),(4,4),(2,2),
\]
de modo que as duas equações são compatíveis. Para $\Re s>0$, $(\partial_\tau+s)$ é invertível nas funções
$4\pi$-periódicas, com inversa $\int_0^\infty e^{-st}a(\tau-t,\xi)\,dt$, e $z=(\partial_\tau+s)^{-1}a_i$ então também
satisfaz $\partial_\xi z=b_i$. Não há obstrução de período, porque $e^{4\pi s}\neq1$. O vínculo
$C(s)h=0$ é necessário para a identidade por trás de $\log W$; as outras compatibilidades decorrem de $L(s)h=0$.

\emph{Geometria.} Com $(q,z,p)$ assim definidos, $\delta$ está na forma RT, tem variáveis RT $e^{s\tau}h$, satisfaz
a condição de Floquet por construção, e resolve as equações linearizadas pelas fórmulas de curvatura do
\cref{lem:R}~(R4), já que todos os $\delta\Omega$ se anulam. Como $h\in\Yp$, os perfis são analíticos em $\tau$ numa
faixa e em $\xi$ dentro da elipse de Bernstein de parâmetro $5/4$, cujo semieixo real é
$\tfrac12(\tfrac54+\tfrac45)=\tfrac{41}{40}$; $q,z,p$ são pares em $\xi$, de modo que $\delta$ é analítico real em $\Mhat$ em
coordenadas cartesianas. A injetividade é clara, pois $\delta$ determina as suas variáveis RT.
\end{proof}

\subsection{Correspondência e multiplicidade física}\label{sec:physical-count}

Combinando os \cref{lem:gauge,lem:R,lem:V}:

\begin{theorem}[Correspondência]\label{thm:correspondence}
Seja $\Re s>0$. A aplicação $\delta\mapsto h$ do \cref{lem:R} induz um isomorfismo linear
\[
  \frac{\{\text{modos de Floquet físicos com expoente }s\}}{\{\text{modos de gauge puro}\}}
  \;\xrightarrow{\ \cong\ }\;
  \frac{\ker L(s)\cap\ker C(s)}{\mathcal G_s},
\]
com $\mathcal G_s$ como no \cref{lem:R}. Os modos físicos generalizados de ordem $k$ correspondem às cadeias de Jordan de
$L$ em $s$ de comprimento $k$ que satisfazem os vínculos de jato
\begin{equation}\label{eq:jet}
  C(s)h_0=0,\qquad C(s)h_j+C_1h_{j-1}=0\quad(1\le j<k),
\end{equation}
módulo $\mathcal G_s$.
\end{theorem}

Note que os vínculos de jato não são as condições $C(s)h_j=0$ para cada $j$; impor estas últimas pode
dar a dimensão errada já em exemplos $2\times2$.

Os vínculos são tratados no espaço de raízes inteiro de uma vez. Fixe uma raiz $s_0$ com $\Re s_0>0$, seja $E$ o seu
espaço de raízes (o autoespaço generalizado de $L(0)$ em $-s_0$, de dimensão finita e contido no domínio em
$\Yp$), e $A=L(0)|_E$. Comparando as potências de $s$ em~\eqref{eq:KCML}, com $K(s)=K(0)+s$, obtém-se $M_1=C_1$,
$M_0=C_0+K(0)C_1-C_1L(0)$ e $K(0)C_0=M_0L(0)$. Defina
\[
  \mathsf D=C_0-C_1L(0)\quad\text{em }E .
\]
Então $K(0)\mathsf D=\mathsf DA$, de modo que $\ker(\mathsf D|_E)$ é $A$-invariante, e para uma cadeia de Jordan
$Ah_j=-s_0h_j-h_{j-1}$ tem-se $\mathsf Dh_j=C(s_0)h_j+C_1h_{j-1}$. Assim, uma cadeia satisfaz os vínculos de
jato~\eqref{eq:jet} se e só se está em $\ker(\mathsf D|_E)$.

\begin{proposition}[Multiplicidade física]\label{prop:physmult}
Seja $s_0$ uma raiz de $L$ com $\Re s_0>0$ e espaço de raízes $E$. A dimensão do espaço dos modos físicos
generalizados com expoente $s_0$, módulo gauge puro, é
\[
  \dim E-\operatorname{rank}(\mathsf D|_E)-\dim(\mathcal G_{s_0}\cap E).
\]
Em particular, se $s_0$ é algebricamente simples com autovetor $h$, esse número é $1$ se $C(s_0)h=0$ e
$h\notin\mathcal G_{s_0}$, e $0$ caso contrário. No cenário de cone fixo da \cref{def:conefixed}, vale o mesmo sem o último
termo.
\end{proposition}

\begin{proof}
Pelo \cref{thm:correspondence}, os modos físicos generalizados módulo gauge correspondem a
$\ker(\mathsf D|_E)/(\mathcal G_{s_0}\cap E)$, e o teorema do núcleo e da imagem dá a fórmula; aqui $\mathcal G_{s_0}\cap E\subset\ker(\mathsf D|_E)$,
porque $\mathsf Dg_*=C(\mu)g_*=0$. O quociente não
exige que $\mathcal G_{s_0}$ tenha um complemento $A$-invariante. Para uma raiz simples, $E=\C h$ e
$\mathsf Dh=C(s_0)h$.

No cenário de cone fixo, seja $\delta$ um modo físico com $\delta g_{++}=0$ no cone. Então
$f_-(\cdot,1)=0$ em (R1), de modo que $y_-(\cdot,1)$ resolve $(\partial_\tau+s-\mu)y_-(\cdot,1)=0$; ela se anula fora da
ressonância, e na ressonância o múltiplo livre de $e^{im_0\tau/2}$ é posto igual a zero. A fixação do gauge
preserva, portanto, o cone e é única, e a obstrução $\kappa$ de (R2) se anula automaticamente. Não
há modos de gauge que preservam o cone com $\Re s>0$: na prova do \cref{lem:gauge}, a condição
$f(0)=0$ exclui a solução constante (\cref{rem:cone-role}). Reciprocamente, os modos do \cref{lem:V} estão na
forma RT, de modo que $\delta g_{++}=0$ em toda parte. A correspondência do \cref{thm:correspondence} vale então com
$\mathcal G_s$ trocado por $0$; para $s=\mu$, o elemento $g_*$ corresponde a $\Lie_{X_0}(\gbar,\phibar)$,
que é um modo físico, mas não um modo de gauge que preserva o cone.
\end{proof}

A \cref{prop:physmult} reduz a parte~(2) do Teorema Principal à parte~(1) e a três fatos sobre os
autovetores: os de $\sB$ e $\sK$ violam os vínculos, o de $\mu$ gera $\mathcal G_\mu$, e o
de $\sphys$ os satisfaz. Eles são provados na \cref{sec:roots}.

\section{A contagem das raízes}\label{sec:counting}

Esta seção prova a parte de contagem do Teorema Principal. Todos os enunciados são para a realização no toro
$\mathbb T(65/64,5/4)$ de $L_A$ e, para $K$, para o toro $\sharp$, $\mathbb T^\sharp(129/128,9/8)$; pelo
\cref{lem:Lreal}, as contagens são as mesmas na realização $\Yp$ de~\cite{RT2019}. Escrevemos
\[
  \Omega_A=(0,3)\times(-\tfrac14,\tfrac14),\qquad \Omega_B=(0,3)\times(\tfrac14,\tfrac34)
\]
para as faixas A e B da \cref{sec:setting-sectors}, truncadas em $\Re s=3$, e $N(L,\Omega)$ para o
número de raízes em $\Omega$ contadas com multiplicidade algébrica.

\begin{theorem}[Contagens]\label{thm:counts}
\begin{enumerate}[label=\textup{(\roman*)},leftmargin=*]
\item $L(s)$ é invertível para $\Re s>2.9261$ e $K(s)$ é invertível para $\Re s>1.765$, para todo $\Im s$.
\item $N(L_A,\Omega_A)=3$ e $N(L_A,\Omega_B)=1$.
\item $L_A$ não tem raízes em $\partial\Omega_A\cup\partial\Omega_B$ além de $s=0$, que é uma raiz
  algebricamente simples.
\end{enumerate}
\end{theorem}

\begin{corollary}\label{cor:four}
Módulo $\tfrac i2\Z$, $L$ tem exatamente quatro raízes em $\{\Re s>0\}$, contadas com multiplicidade algébrica. No
eixo imaginário a sua única raiz é $s\equiv0$, que é algebricamente simples.
\end{corollary}

\begin{proof}
Pelo \cref{lem:sectorB}, as raízes de $L$ módulo $\tfrac i2\Z$ correspondem, com multiplicidades, às raízes de
$L_A$ módulo $i\Z$, que são contadas na faixa $\widetilde Q=\{-\tfrac14<\Im s\le\tfrac34\}$. Por~(i),
não há nenhuma com $\Re s\ge3$, e por~(iii) nenhuma nas retas $\Im s=\tfrac14$ e $\Im s=-\tfrac14\equiv\tfrac34$
com $0<\Re s<3$. Assim, as raízes com $\Re s>0$ são as de $\Omega_A\cup\Omega_B$, e são $3+1$. No eixo, o
segmento $\{\Re s=0,\ -\tfrac14\le\Im s\le\tfrac34\}$ é um período completo de $L_A$ e está contido
em $\partial\Omega_A\cup\partial\Omega_B$, de modo que~(iii) se aplica.
\end{proof}

A prova do \cref{thm:counts}~(i) está na \cref{sec:counting-large}. As partes (ii) e (iii) são provadas nas
\cref{sec:counting-rouche,sec:counting-reference,sec:counting-windows,sec:counting-finite} por uma
versão para operadores do teorema de Rouch\'e: após uma deflação de posto dois das raízes $0$ e $\mu$, $L_A$ é
comparado em $\partial\Omega$ com uma referência cujos zeros são os autovalores de uma matriz $14016\times14016$,
e estes são contados a partir de uma decomposição de Schur com resíduo certificado.

\subsection{Parte real grande}\label{sec:counting-large}

No espaço do toro, os operadores do fundo são pontuais. Com $w=e^{i\tau/2}$ e $z=e^{i\theta}$ no
toro $|w|=a$, $|z|=b$, a multiplicação por um campo $f$ vira a multiplicação pela sua extensão analítica
$f(w,z)$, a reflexão $P$ vira $z\mapsto-z$, e a divisão por $\xi$ vira a multiplicação por
$1/\xi(z)$, $\xi(z)=(z+z^{-1})/2$, que é limitada, pois $b>1$. Consequentemente,
$\Re\langle -Bx,x\rangle$ é a integral sobre o toro de uma forma hermitiana nos sete valores
$(x_1(z),x_2(\pm z),x_3(\pm z),x_4(\pm z))$ (usando $x_1(-z)=-x_1(z)$), com coeficientes dados pelo fundo,
e
\[
  \max\Re W(-B)\le\sup_{\text{toro}}\lambda_{\max}\big(N^{-1/2}\operatorname{Herm}(-\mathcal S)N^{-1/2}\big),
\]
onde $\mathcal S$ é o símbolo $7\times7$ e $N$ leva em conta a multiplicidade dos valores. Aqui
$W(T)=\{\langle Tx,x\rangle:\ x\in\Dom(T),\ \|x\|=1\}$ é a imagem numérica, onde $\Dom(T)$ é o domínio
de $T$. Essa cota captura todos os cancelamentos
entre componentes e fases de Fourier.

\begin{lemma}[Imagem numérica do fundo]\label{lem:F3}
Em $\mathbb T(65/64,5/4)$, $\max\Re W(-B)\le2.9492$ e $\|B\|\le3.9642$; para a parte livre,
$\max\Re W(-J)\le-0.02319$. Em $\mathbb T^\sharp(129/128,9/8)$, $\max\Re W(-B^\sharp)\le1.4396$ e
$\|B^\sharp\|\le1.6497$, e $\max\Re W(-K(0))\le 1.765$.
\end{lemma}

\begin{proof}[Método da prova]
O supremo de $\lambda_{\max}$ do símbolo sobre o toro é cotado numa grade de $2048\times4096$
centros para $L$ e de $1024\times2048$ para $K$. Em cada centro, o símbolo é avaliado em aritmética de bolas (Arb), e a desigualdade é verificada por uma
fatoração $LDL^{\mathsf T}$ em bolas da matriz deslocada. Entre os centros, constantes de Lipschitz dos
campos do fundo, calculadas em aritmética racional exata, dão a folga, também avaliada em aritmética de bolas.
A folga de Lipschitz somada entre os centros é no máximo $0.0772$ para $L$ e $0.0536$ para $K$. A parte livre é tratada por uma forma
fechada. Os detalhes e as referências aos arquivos estão na \cref{sec:computation}.
\end{proof}

\begin{proof}[Prova do \cref{thm:counts}~(i)]
Para $x$ no domínio, $\Re\langle L(s)x,x\rangle\ge(\Re s-2.9492+0.02319)\|x\|^2$, logo
$\|L(s)x\|\ge(\Re s-2.9261)\|x\|$, e $L(s)$ é injetivo para $\Re s>2.9261$. Como $L(s)$ é de Fredholm de
índice zero (\cref{prop:fredholm}), ele é invertível. O argumento para $K$ é o mesmo.
\end{proof}

A cota $2.9261$ deve ser comparada com a maior raiz, $\sphys\approx0.733$: a exclusão não é ótima,
e é por isso que os retângulos $\Omega_A,\Omega_B$ vão até $\Re s=3$. Antes do uso do símbolo, majorantes bloco a
bloco davam uma cota duas ordens de grandeza maior.

\subsection{O teorema de Rouch\'e para famílias de operadores, e a deflação}\label{sec:counting-rouche}

Usamos a versão para operadores do teorema de Rouch\'e de Gohberg e Sigal~\cite{GohbergSigal1971}, na sua
forma de homotopia. Seja $\Omega$ um domínio limitado com fronteira suave por partes, e $H,R$ famílias analíticas de
operadores de Fredholm de índice zero numa vizinhança de $\overline\Omega$. Suponha ainda que $H-R$ seja compacto. Se
$R+t(H-R)$ é invertível em $\partial\Omega$ para todo $t\in[0,1]$, então $H$ e $R$ têm o mesmo número de zeros em $\Omega$, contados com
multiplicidade algébrica. Uma condição suficiente é $\|R(s)^{-1}(H(s)-R(s))\|<1$ em $\partial\Omega$.

\emph{Normalização.} A parte livre $J+s$ é triangular no índice de Chebyshev dentro de cada modo de Fourier, com
diagonal $\mu(n+1)+s+im/2$, de modo que $Q_0(s):=(J+s)^{-1}$ é analítico e invertível em $\Re s>-\mu$. A família
$H(s)=\widetilde L(s)Q_0(s)$ tem os mesmos zeros que $\widetilde L$, com as mesmas multiplicidades, e
$H(s)=I+(B+E)Q_0(s)$, com $E=\widetilde L-L$ a deflação abaixo, é a identidade mais um operador compacto.

\emph{Deflação.} As raízes $0$ e $\mu$ são conhecidas exatamente: $L(s)x_0=s\,x_0$ para $x_0=\partial_\tau\omegabar$
(invariância por translação em $\tau$), e $L(s)x_1=(s-\mu)\,x_1$ para $x_1=g_*$ (\cref{sec:physical-defs}). A
raiz $0$ está em $\partial\Omega_A$, e a raiz $\mu$ mantém o resolvente grande entre $0$ e $\mu$. Removemos
ambas por uma perturbação de posto finito.

\begin{lemma}[Deflação de Brauer]\label{lem:brauer}
Sejam $\phi_0,\phi_1$ funcionais limitados e $\delta_0,\delta_1\in\C$, e ponha
$\widetilde L(s)=L(s)+\delta_0x_0\phi_0^*+\delta_1x_1\phi_1^*$. Então
\[
  \det\big(I+L(s)^{-1}(\widetilde L(s)-L(s))\big)=\frac{q(s)}{s(s-\mu)},
\]
onde $q$ é o polinômio quadrático $\det\big(\operatorname{diag}(s,s-\mu)+[\delta_j\phi_i^*x_j]_{ij}\big)$, e
para todo $s_0$ as multiplicidades algébricas satisfazem
\[
  \mult_{\widetilde L}(s_0)=\mult_L(s_0)+\operatorname{ord}_{s_0}\frac{q(s)}{s(s-\mu)} .
\]
\end{lemma}

\begin{proof}
Como $L(s)^{-1}x_0=x_0/s$ e $L(s)^{-1}x_1=x_1/(s-\mu)$, a perturbação $L(s)^{-1}(\widetilde L-L)$ tem
posto dois com imagem em $\operatorname{span}\{x_0,x_1\}$, e o determinante se reduz a um determinante
$2\times2$. A fórmula das multiplicidades é a forma de resíduo logarítmico da teoria de Gohberg--Sigal para a
fatoração $\widetilde L=L\,(I+L^{-1}(\widetilde L-L))$, na qual o segundo fator é uma perturbação meromorfa de
posto finito da identidade cuja multiplicidade logarítmica é a ordem do seu determinante.
\end{proof}

Tomamos $\phi_i$ com suporte na caixa finita $Z$ abaixo, biortogonais a aproximações explícitas $x_j^A$ calculadas
a partir dos dados $\omega_A$ ($\phi_i^*x_j^A=\delta_{ij}$), e $(\delta_0,\delta_1)=(1,1+\mu_A)$. Com
$e_{ij}=\phi_i^*(x_j-x_j^A)$,
\[
  q(s)=(s+1+e_{00})\big(s-\mu+\delta_1(1+e_{11})\big)-\delta_0\delta_1e_{01}e_{10},
\]
de modo que as duas raízes deflacionadas se movem para perto de $-1$. Os erros satisfazem $\|P_Z(x_0-x_0^A)\|\le1.5\cdot10^{-9}$
e $\|x_1-x_1^A\|\le2.51\cdot10^{-8}$, logo $|e_{ij}|\le\|\phi_i\|\,\|x_j-x_j^A\|\le5.51\cdot2.51\cdot10^{-8}$, e os dois fatores de $q$ têm módulo
$\ge1-10^{-6}$ em $\Re s\ge0$. Portanto $q$ não tem zeros em $\Re s\ge0$, e pelo \cref{lem:brauer}
\begin{equation}\label{eq:brauer-count}
  N(L_A,\Omega_A)=N(\widetilde L_A,\Omega_A)+1,\qquad N(L_A,\Omega_B)=N(\widetilde L_A,\Omega_B),
\end{equation}
o $+1$ vindo do polo de $q(s)/(s(s-\mu))$ em $\mu\in\Omega_A$. Além disso, em todo $s$ com $\Re s\ge0$ no qual
$\widetilde L_A(s)$ é invertível, $\mult_L(s)=1$ se $s\in\{0,\mu\}$ e $\mult_L(s)=0$ caso contrário. Nos certificados, a
deflação usa $x_j^A$, e as diferenças $x_j-x_j^A$ entram como perturbações.

\subsection{A referência e a sua verificação no contorno}\label{sec:counting-reference}

Particione o espaço de coeficientes do setor A em
\begin{itemize}[leftmargin=*]
\item a caixa finita $Z=\{|m|\le31,\ 0\le n<128\}$, de dimensão $14016$;
\item $26$ \emph{janelas} (\emph{windows}) $J$ de no máximo $16$ modos de Fourier consecutivos, que cobrem o resto da caixa
  $\{|m|<200,\ n<400\}$;
\item a parte \emph{distante} (\emph{far part}), o complemento.
\end{itemize}
Sejam $P_Z,P_J,P_{\mathrm{far}}$ as projeções coordenadas e $\widehat H_{ab}(s)=P_aH(s)P_b$. A referência é a
família bloco-diagonal
\[
  R(s)=\operatorname{diag}\big(\widehat H_{ZZ}(s),\ \widehat H_{JJ}(s)\ (J\text{ janelas}),\ I_{\mathrm{far}}\big).
\]
Ela é analítica em $\Re s>-\mu$, e os seus zeros em $\Omega$ são os dos blocos. Nos certificados, o bloco
finito é tomado um passo mais longe do operador: com $W$ a matriz diagonal dos pesos, $\widehat H_{ZZ}$ é
trocado por $R_Z(s)=W^{-1}(\mathsf A+s)W\,(P_Z(J+s)P_Z)^{-1}$, onde $\mathsf A$ é a matriz gravada, em ponto flutuante,
que aproxima $W\widetilde L_Z(0)W^{-1}$, construída a partir dos dados $\omega_A$ e dos vetores $x_j^A$ e tomada como
matriz exata. A diferença $P_Z\widetilde LP_Z-W^{-1}\mathsf AW$, devida ao erro do fundo, aos vetores da deflação e
ao arredondamento na montagem, é parte de $H-R$ e é cotada nas entradas abaixo.

\begin{lemma}[O bloco finito]\label{lem:ZZ}
$\widehat H_{ZZ}(s)=\widetilde L_Z(s)\,(P_Z(J+s)P_Z)^{-1}$, onde $\widetilde L_Z(s)=P_Z\widetilde L(s)P_Z$ é o
truncamento $14016\times14016$. Em particular, os zeros de $\widehat H_{ZZ}$ em $\Omega$, com multiplicidades, são
os pontos $s$ com $-s$ autovalor da matriz $\widetilde L_Z(0)$, contados com multiplicidade algébrica. O mesmo
vale para $R_Z$, com $\mathsf A$ no lugar de $W\widetilde L_Z(0)W^{-1}$.
\end{lemma}

\begin{proof}
$J+s$ preserva o índice de Fourier e é triangular superior em $n$, de modo que a imagem de $P_Z$ é invariante por
$J+s$ e por $Q_0(s)$: $Q_0(s)P_Z=P_ZQ_0(s)P_Z=(P_Z(J+s)P_Z)^{-1}$. Portanto
$\widehat H_{ZZ}(s)=P_Z\widetilde L(s)Q_0(s)P_Z=P_Z\widetilde L(s)P_Z\,(P_Z(J+s)P_Z)^{-1}$, e o segundo fator é
invertível e analítico em $\Re s>-\mu$.
\end{proof}

\emph{O critério de Perron.} Para verificar que $R+t(H-R)$ é invertível em $\partial\Omega$ para todo $t\in[0,1]$,
usamos inversas aproximadas dos blocos diagonais, $V_Z(s)\approx\widehat H_{ZZ}(s)^{-1}$, $V_J\approx\widehat
H_{JJ}(c)^{-1}$ num centro fixo $c$, e $I$ para a parte distante, e cotamos todas as entradas dos blocos.

\begin{lemma}[Critério de Perron]\label{lem:perron}
Sejam $X=\bigoplus_iX_i$, $A=R+E$ com $R=\operatorname{diag}(R_i)$, e $V_i$ operadores limitados com
$\|I-V_iR_i\|\le\rho_i$ e $\|V_iE_{ij}\|\le n_{ij}$. Se o raio espectral $\theta$ da matriz não negativa
$\operatorname{diag}(\rho_i)+(n_{ij})$ é $<1$, então $R+tE$ é injetivo para todo $t\in[0,1]$; se $R+tE$ é de Fredholm de
índice zero, ele é invertível.
\end{lemma}

\begin{proof}
Com $V=\operatorname{diag}(V_i)$, $V(R+tE)=I-\big((I-VR)-tVE\big)$. Na norma
$\max_iw_i^{-1}\|x_i\|$, com $w$ um vetor de Perron positivo de uma matriz ligeiramente maior, o operador entre
parênteses tem norma $\le\theta+\epsilon<1$, já que as normas dos seus blocos são cotadas entrada a entrada pela matriz acima
(e são crescentes em $t$). Assim $V(R+tE)$ é invertível, e $R+tE$ é injetivo.
\end{proof}

Os níveis são $Z$, as janelas e a parte distante; esta é a cota de \emph{Perron em três níveis}. As entradas são
calculadas como segue (\cref{sec:computation}).
\begin{itemize}[leftmargin=*]
\item \emph{Nível $Z$.} $V_Z(s)$ é aplicado por meio de uma decomposição de Schur de $\mathsf A$, usada
  apenas como resolvedor; os resíduos são certificados a posteriori. Uma cota $t_p\ge\|(\mathsf A+p)^{-1}\|$ é certificada por uma desigualdade na ordem de Loewner, $t^2MM^*-I\succeq0$ (com precondicionamento perto da
  raiz da faixa B). O acoplamento $Z\leftarrow$ cauda é fatorado como uma parte de posto baixo $F\Phi$ com um resto
  certificado; com a escala $D_Z=J_Z(c)Q_{ZZ}(s)$, a entrada cauda$\,\leftarrow Z$ fica independente de $s$.
\item \emph{Janelas.} Cada $V_J$ é a inversa exata do bloco calculado no centro, e os blocos são
  cotados por estimativas de Jacobi por blocos; a dependência em $s$ é cotada por $\rho_J(s)\le\rho_J(c)+d\,K_J/(1-d\,n_{Q,J})$,
  $d=|s-c|$, com constantes calculadas no centro (e uma fórmula refinada para as janelas centrais).
\item \emph{Parte distante.} Formas fechadas para $\|Q_0(s)P_{\mathrm{far}}\|$, uniformes em discos.
\item \emph{Erros do fundo e da deflação.} A diferença entre $\omegabar$ e os dados, e os erros
  $x_j-x_j^A$, entram como perturbações globais de todas as entradas.
\end{itemize}

\emph{Discos.} A cota do \cref{lem:perron} é tornada uniforme num disco em torno de cada ponto $p$ do contorno:
a variação das entradas do nível $Z$ é controlada em primeira ordem por quantidades estruturadas calculadas com
resoluções de Schur, e em segunda ordem pela cota do resolvente, usando que $\sigma_{\min}(\widetilde L_Z+s)$ é
$1$-Lipschitz em $s$. Para cada disco, o raio é escolhido por bissecção sobre uma estimativa de $\theta$ em ponto flutuante, com alvo $0.999$; as
cotas racionais certificadas são então calculadas para o raio escolhido, e podem passar ligeiramente do alvo, desde que
fiquem abaixo de $1$. O contorno de $\Omega_A$ é
coberto por $96$ discos, com $\max\theta\le0.99842$, e o de $\Omega_B$ pelos $96$ discos de $\Omega_A$ na sua
aresta inferior e mais $78$ discos, com $\max\theta\le0.99908$; os raios vão de $0.0028$ perto da raiz da faixa B
a $0.4995$ em $\Re s\ge2$. A cobertura das quatro arestas pela união dos discos é conferida em aritmética
racional. Isso prova que $R+t(H-R)$ é invertível em $\partial\Omega_A$ e em $\partial\Omega_B$, logo
\[
  N(\widetilde L_A,\Omega)=N(R,\Omega)\qquad(\Omega=\Omega_A,\Omega_B),
\]
e que $\widetilde L_A$ é invertível em $\partial\Omega_A\cup\partial\Omega_B$. Pela observação após
\eqref{eq:brauer-count}, isso é o \cref{thm:counts}~(iii).

\subsection{As janelas não têm zeros}\label{sec:counting-windows}

O critério de Perron no contorno dá $N(\widetilde L,\Omega)=N(R,\Omega)=N(\widehat H_{ZZ},\Omega)+
\sum_JN(\widehat H_{JJ},\Omega)$, mas não diz nada sobre zeros dos blocos das janelas dentro de $\Omega$.

\begin{lemma}\label{lem:windows}
Para toda janela $J$, $\widehat H_{JJ}(s)$ é invertível para todo $s\in\overline{\Omega_A}\cup\overline{\Omega_B}$.
\end{lemma}

\begin{proof}
Particione cada faixa em $10$ retângulos $\Omega_k$, um para cada centro de cauda $c_k$. Em $\Omega_k$, a função
$F_k(s)=I-V_J(c_k)\widehat H_{JJ}(s)$ é analítica, de modo que $\|F_k(s)\|$ é subharmônica e atinge o seu máximo em
$\partial\Omega_k$. Na fronteira, $\|F_k\|$ é cotada pelas mesmas estimativas de sensibilidade dos discos,
avaliadas em discos em torno de pontos de amostra da fronteira. O máximo é $0.1488$ na faixa A e $0.1517$ na faixa B.
Assim $\|F_k\|<1$ em $\Omega_k$, e $\widehat H_{JJ}(s)$ é invertível.
\end{proof}

\subsection{A contagem finita}\label{sec:counting-finite}

Resta contar os autovalores da matriz gravada $\mathsf A$ em $-\Omega$. Uma decomposição de Schur em ponto flutuante $\mathsf A\approx UTU^*$ dá a matriz $\mathsf B=UTU^{-1}$,
cujos autovalores são exatamente as entradas diagonais de $T$ (com $U$ invertível).

\begin{lemma}[Contagem finita]\label{lem:schur}
Sejam $\varepsilon_B\ge\|\mathsf A-\mathsf B\|$ e $t(s)\ge\|(\mathsf A+s)^{-1}\|$ em $\partial\Omega$. Se
$\sup_{\partial\Omega}t(s)\,\varepsilon_B<1$, então $\mathsf A$ tem exatamente $\#\{i:-T_{ii}\in\Omega\}$ autovalores
em $-\Omega$, com multiplicidade.
\end{lemma}

\begin{proof}
$\mathsf A+s+t'(\mathsf B-\mathsf A)$ é invertível em $\partial\Omega$ para $t'\in[0,1]$, por uma série de Neumann, de modo que o
número de autovalores no interior não muda ao longo da homotopia.
\end{proof}

Temos $\varepsilon_B\le\|\mathsf AU-UT\|_F/\sqrt{1-\|I-U^*U\|_F}$ com $\|\mathsf AU-UT\|_F\le5.97\cdot10^{-5}$ e
$\|I-U^*U\|_F\le1.31\cdot10^{-6}$, e $t(s)\le t_p/(1-rt_p)$ em cada disco, de modo que $\sup t\le559.2$ e
$t\,\varepsilon_B\le0.0334$. A diagonal de $T$ tem exatamente duas entradas $-T_{ii}$ em $\Omega_A$, em
$0.401025\ldots$ e $0.733180\ldots$, e exatamente uma em $\Omega_B$, em $0.0414194\ldots+0.5i$; a entrada mais próxima
fora de $\Omega_A$ está à distância $0.168$ da sua fronteira. Portanto $N(\widehat H_{ZZ},\Omega_A)=2$ e
$N(\widehat H_{ZZ},\Omega_B)=1$.

\begin{proof}[Prova do \cref{thm:counts}~(ii)]
Pelos \cref{lem:windows,lem:schur,lem:ZZ}, $N(R,\Omega_A)=2$ e $N(R,\Omega_B)=1$. Pela verificação de Perron nos
contornos, $N(\widetilde L_A,\Omega)=N(R,\Omega)$ e, por~\eqref{eq:brauer-count},
$N(L_A,\Omega_A)=2+1=3$ e $N(L_A,\Omega_B)=1$.
\end{proof}

As posições $0.401$, $0.733$ e $0.0414+0.5i$ da referência finita não são usadas na contagem; as raízes de
$L$ são localizadas de modo independente na \cref{sec:roots}. O argumento de Rouch\'e garante apenas que há três
raízes em $\Omega_A$ e uma em $\Omega_B$.

\section{Identificação das raízes}\label{sec:roots}

Pelo \cref{thm:counts}, $L_A$ tem três raízes em $\Omega_A$ e uma em $\Omega_B$. Esta seção localiza cada uma
delas, prova que cada uma é algebricamente simples e determina o comportamento do seu autovetor em relação
aos vínculos. Junto com a \cref{prop:physmult}, isso prova o Teorema Principal (\cref{sec:roots-proof}).

\subsection{Newton--Kantorovich para o sistema aumentado}\label{sec:roots-nk}

Fixe $c_0$ com $\Re c_0>-\mu$, seja $Q_0=(J+c_0)^{-1}$, e seja $\ell$ um funcional limitado. Os autopares
$(h,\lambda)$ de $L$ com $\ell(h)=1$ correspondem, por meio de $h=Q_0y$, aos zeros da aplicação do sistema aumentado
\[
  F(y,\lambda)=\big(L(\lambda)Q_0\,y,\ \ell(Q_0y)-1\big).
\]
Como $L(\lambda)=L(0)+\lambda$, $F$ é quadrática, com
$DF(y,\lambda)[\eta,\nu]=\big(L(\lambda)Q_0\eta+\nu\,Q_0y,\ \ell(Q_0\eta)\big)$ e derivada segunda constante.

\begin{lemma}[Simplicidade a partir do sistema aumentado]\label{lem:bordered}
Se $F(y_*,\lambda_*)=0$, então $DF(y_*,\lambda_*)$ é invertível se e só se $\lambda_*$ é uma raiz algebricamente
simples de $L$.
\end{lemma}

\begin{proof}
Seja $h_*=Q_0y_*$. Se $DF[\eta,\nu]=0$ com $\nu\neq0$, então $k=Q_0\eta/\nu$ satisfaz $L(\lambda_*)k=-h_*$,
uma cadeia de Jordan de comprimento dois. Se $\nu=0$, então $Q_0\eta\in\ker L(\lambda_*)$ e $\ell(Q_0\eta)=0$; se o
núcleo é gerado por $h_*$, com $\ell(h_*)=1$, isso obriga $\eta=0$. Portanto $DF$ é injetivo se e só se
$\ker L(\lambda_*)=\C h_*$ e não há cadeia de Jordan, isto é, se e só se $\lambda_*$ é algebricamente
simples. $DF$ é uma perturbação de posto finito de um operador de Fredholm de índice zero, de modo que a injetividade
equivale à invertibilidade.
\end{proof}

Usamos o teorema de Newton--Kantorovich na seguinte forma padrão~\cite{NakaoPlumWatanabe2019}. Sejam
$x_0=(y_0,\lambda_0)$ um zero aproximado e $\mathcal A$ uma inversa aproximada de $DF(x_0)$, e suponha que, em
alguma norma,
\[
  \|I-\mathcal A\,DF(x_0)\|\le\theta,\qquad \tfrac12\|\mathcal A\,D^2F\|\le Z_2,\qquad \|\mathcal A F(x_0)\|\le Y.
\]
Se $\mathcal A$ é injetivo (aqui ele é invertível por construção) e $(1-\theta)^2>4Z_2Y$, então $T(x)=x-\mathcal AF(x)$ é uma contração da bola $B(x_0,r)$,
$r=\big(1-\theta-\sqrt{(1-\theta)^2-4Z_2Y}\big)/(2Z_2)$, de modo que $F$ tem um único zero $x_*$ ali, e
$\|I-\mathcal A\,DF(x_*)\|<1$, de modo que $DF(x_*)$ é invertível. A unicidade vale em toda bola $B(x_0,\rho)$ com
$\theta+2Z_2\rho<1$.

A norma é um máximo ponderado de normas de blocos sobre os mesmos três níveis da \cref{sec:counting-reference}:
a caixa finita aumentada $Z'$ ($Z$ mais a linha e a coluna acrescentadas), as $26$ janelas e a parte distante. A
cota $\theta$ é a raiz de Perron da matriz das cotas dos blocos, verificada em aritmética racional, e $\mathcal A$
é construído a partir da inversa exata da matriz aumentada calculada em $Z'$ e dos blocos calculados das janelas. O
centro $\lambda_0$ é um autovalor do truncamento $32\times128$ refinado pela iteração de Newton no sistema aumentado, e
$h_0$ é obtido por iteração inversa na matriz ponderada.

\begin{proposition}[Localização e simplicidade]\label{prop:nk}
Para cada linha da tabela, existe um único zero $x_*=(y_*,\lambda_*)$ de $F$ em $B(x_0,r)$, e $\lambda_*$ é uma
raiz algebricamente simples de $L_A$ com
\begin{center}
\begin{tabular}{@{}lllllll@{}}
\toprule
raiz & $\lambda_0$ & $\theta$ & $Z_2$ & $Y$ & $r$ & $|\lambda_*-\lambda_0|\le$ \\
\midrule
$\mu$ & $\mu_A=0.16830707\ldots$ & $0.880172$ & $183.66$ & $6.60\cdot10^{-6}$ & $6.07\cdot10^{-5}$ & $2.82\cdot10^{-5}$ \\
$\sK$ & $0.401024733$ & $0.839937$ & $42.97$ & $5.45\cdot10^{-5}$ & $3.79\cdot10^{-4}$ & $1.72\cdot10^{-4}$ \\
$\sphys$ & $0.7331796596606924$ & $0.805721$ & $10.33$ & $4.19\cdot10^{-6}$ & $2.16\cdot10^{-5}$ & $1.0163\cdot10^{-5}$ \\
$\sB+\frac i2$ & $0.0414194105346554+0.5i$ & $0.91153$ & $155.02$ & $5.97\cdot10^{-6}$ & $7.81\cdot10^{-5}$ & $3.67\cdot10^{-5}$ \\
\bottomrule
\end{tabular}
\end{center}
Todas as quantidades são arredondadas para fora a partir dos valores racionais exatos dos arquivos de certificado: para cima no caso de $\theta$,
$Z_2$, $Y$, $r$ e do erro de $\lambda$. Os valores de $Z_2$ cotam o próprio $\|\mathcal A\,D^2F\|$, o que é
conservador por um fator dois. As raízes das três primeiras linhas são as três raízes de $L_A$ em $\Omega_A$, e a última é a raiz em
$\Omega_B$.
\end{proposition}

\begin{proof}
As condições de Newton--Kantorovich são verificadas com as quantidades da tabela (\cref{sec:computation}), e o
\cref{lem:bordered} dá a simplicidade. Os intervalos das três primeiras linhas são dois a dois disjuntos e contidos
em $\Omega_A$, e cada um contém uma raiz simples; como $N(L_A,\Omega_A)=3$, essas são todas as raízes em $\Omega_A$.
O mesmo vale para $\Omega_B$, onde $N(L_A,\Omega_B)=1$.
\end{proof}

\subsection{A raiz de gauge}\label{sec:roots-gauge}

\begin{proposition}[A raiz de gauge]\label{prop:mu}
$\mu$ é uma raiz algebricamente simples de $L_A$, e $\ker L_A(\mu)=\C\,g_*$, com $g_*=(1+Z)\omegabar$.
\end{proposition}

\begin{proof}
$L(\mu)g_*=0$ exatamente (\cref{sec:physical-defs}), e $g_*$ está no domínio, já que $g_*-g_A$ é pequeno em norma,
onde $g_A=(1+Z)\omega_A$ é calculado a partir dos dados, e $J+\lambda_0$ é injetivo no domínio maximal. Normalize
$g_n=g_*/\ell(g_*)$ e seja $x_g=(Q_0^{-1}g_n,\mu)$, um zero exato de $F$. Mostramos que $x_g\in B(x_0,r)$ para a primeira
linha da \cref{prop:nk}; então $x_g=x_*$ pela unicidade, de modo que a raiz simples encontrada ali é exatamente $\mu$, com
autovetor $g_*$. Pela equação de autovalor,
\[
  y_g-y_0=-B(g_n-h_0)-L(\lambda_0)h_0+(\lambda_0-\mu)g_n-\delta\mu\,J_1(g_n-h_0),
\]
onde $\delta\mu=\mu-\mu_A$ e $J_1$ é a derivada de $J$ em relação a $\mu$. Os termos são cotados por
$\|g_*-g_A\|\le2.51\cdot10^{-8}$, que usa as cotas publicadas de~\cite{RT2019} para a diferença entre
$\omegabar$ e os dados, junto com uma cota de $(1+Z)$ aplicado à correção, e pelo resíduo calculado
$\|L(\lambda_0)h_0\|\le3.76\cdot10^{-6}$. O resultado é $\|x_g-x_0\|\le1.82\cdot10^{-5}$, dentro da bola $B(x_0,\rho)$, $\rho=3.26\cdot10^{-4}$, na qual
$\theta+2Z_2\rho<1$ e o zero de $F$ é único.
\end{proof}

\subsection{Raízes que violam os vínculos}\label{sec:roots-witness}

\begin{proposition}[Testemunhos dos vínculos, \emph{constraint witnesses}]\label{prop:witness}
Seja $h_*=Q_0y_*$ o autovetor da \cref{prop:nk} em $\sK$, respectivamente em $\sB+\tfrac i2$. Então, em
$\mathbb T^\sharp(129/128,9/8)$,
\[
  \|\Pi_F\,C(\sK)h_*\|\ge0.16975,\qquad \|\Pi_F\,C(\sB+\tfrac i2)h_*\|\ge0.35469,
\]
onde $\Pi_F$ é a projeção na caixa $F=\{|m|<12,\ n<48\}$. Em particular, os dois autovetores violam os
vínculos, e o mesmo acontece com o autovetor de $L$ em $\sB$ (\cref{lem:sectorB}).
\end{proposition}

\begin{proof}
No centro, $c_{\mathrm{ref}}=\|\Pi_FC(\lambda_0)h_0\|$ é calculado rigorosamente na caixa. A diferença
$\Pi_F\big(C(\lambda_*)h_*-C(\lambda_0)h_0\big)$ é cotada nível a nível usando a bola de Newton--Kantorovich:
o nível $Z'$ por meio de $\|\Pi_FCQ_0P_Z\|$, as janelas pelo decaimento das colunas (as colunas da cauda começam em
$n=128$ e só chegam à caixa através do fator $(5/4)^{-80}$), a parte distante por formas fechadas, e a variação
de $\lambda$ por meio de $C_1$. Para $\sK$: $c_{\mathrm{ref}}\ge0.171739$, e as perdas somam no máximo $1.981\cdot10^{-3}$, o que dá
$0.171739-0.001981=0.169758$. Para $\sB$: $c_{\mathrm{ref}}\ge0.355188$, e as perdas somam no máximo
$4.954\cdot10^{-4}$, o que dá $0.355188-0.0004954=0.3546926$.
\end{proof}

\subsection{A raiz física: injetividade de $K$}\label{sec:roots-K}

Não se pode mostrar por uma cota inferior que o autovetor em $\sphys$ satisfaz os vínculos; em vez disso, mostramos que
ele não pode violá-los. Por~\eqref{eq:KCML}, se $L(s)h=0$ então $K(s)C(s)h=0$, de modo que $C(s)h\neq0$ só é possível numa
raiz de $K$.

\begin{proposition}[Injetividade de $K$]\label{prop:Kinj}
Em $\mathbb T^\sharp(129/128,9/8)$, $K(s)$ é injetivo para todo $s$ com $\Re s\ge0$ e $0\le\Im s\le\tfrac14$,
exceto num ponto $s=\sK'$, com $|\sK'-0.401|<0.05$, que é uma raiz algebricamente simples de $K$.
\end{proposition}

\begin{proof}
A região é coberta por três certificados.
\begin{itemize}[leftmargin=*]
\item $\Re s>1.765$: \cref{thm:counts}~(i).
\item $\{0\le\Re s\le1.765,\ 0\le\Im s\le\tfrac14\}$ menos o disco aberto $|s-0.401|<0.05$: uma cobertura por $154$
  células (\emph{tiles}). Em cada célula, a injetividade de $K(s)$ decorre do \cref{lem:perron} aplicado a $K(s)Q_0(c)$, com quatro
  níveis (uma caixa finita, duas camadas e a parte distante) e raiz de Perron $\le0.99612$; a cobertura é conferida em
  aritmética racional por uma árvore quaternária (\emph{quadtree}).
\item O disco $|s-0.401|\le0.05$: um argumento de Rouch\'e para operadores no seu círculo de fronteira, como na
  \cref{sec:counting-reference}, com referência $\operatorname{diag}(\widehat H_{ZZ}(s),I,I,I)$ numa caixa
  $Z=20\times80$ e raízes de Perron $\le0.8126$ em $16$ células que cobrem o círculo (a norma com sinal não é simétrica
  sob conjugação, de modo que o círculo inteiro é coberto). A referência finita é contada por um argumento de Rouch\'e
  escalar para o sistema aumentado: o truncamento tem exatamente um autovalor no disco, simples, perto de $0.4010247311$.
\end{itemize}
\end{proof}

\begin{corollary}\label{cor:sK}
$\sK'=\sK$, e o autovetor de $L$ em $\sphys$ satisfaz $C(\sphys)h=0$.
\end{corollary}

\begin{proof}
Os autovetores da \cref{prop:nk} são calculados no espaço do toro; como $|\sK|,|\sphys|\le3.1$, eles estão no
domínio de $L$ em $\Yp$ pelo \cref{lem:Lreal}, de modo que o \cref{lem:KCML} se aplica a eles com $(a,b)=(65/64,5/4)$ e
$(a',b')=(129/128,9/8)$. Pela \cref{prop:witness}, $C(\sK)h_*\neq0$, e $K(\sK)C(\sK)h_*=M(\sK)L(\sK)h_*=0$, de modo que $\sK$ é
uma raiz de $K$ na região da \cref{prop:Kinj}; portanto $\sK=\sK'$. Como $\sphys$ é real e $|\sphys-0.401|>0.33$,
$K(\sphys)$ é injetivo, de modo que $C(\sphys)h=0$ para o autovetor $h$. Aqui $C(\sphys)h$ está no domínio de $K$ em $Y^\sharp(129/128,9/8)\subset\mathbb T^\sharp(129/128,9/8)$, onde
a \cref{prop:Kinj} se aplica (\cref{app:KCML}).
\end{proof}

\subsection{Realidade}\label{sec:roots-reality}

\begin{proposition}\label{prop:real}
$\sK$ e $\sphys$ são reais, e $\Im(\sB+\tfrac i2)=\tfrac12$, isto é, a raiz $\sB$ de $L$ é real.
\end{proposition}

\begin{proof}
Na realização $\Yp$, a conjugação $h_{mn}\mapsto\overline{h_{-m,-n}}$ é uma isometria, e
$\overline{L_A(\bar s)}=L_A(s)$ no sentido da \cref{sec:setting-sectors}; assim, as raízes de $L_A$ em $\Yp$, com
multiplicidades, são simétricas por $s\mapsto\bar s$, e as da faixa B por $s\mapsto\bar s+i$. Pelo
\cref{lem:Lreal}, o mesmo vale na realização do toro para $|s|\le3.1$. Em $\Omega_A$ as raízes são $\mu,\sK,\sphys$,
com $\sK$ e $\sphys$ em discos de raio $<2\cdot10^{-4}$ em torno de centros reais; a conjugada de $\sK$ é uma raiz no
mesmo disco, logo igual a $\sK$, e o mesmo para $\sphys$. Em $\Omega_B$ há uma única raiz, de modo que ela é fixada por
$s\mapsto\bar s+i$, cujo conjunto de pontos fixos é $\Im s=\tfrac12$.
\end{proof}

\subsection{Prova do Teorema Principal}\label{sec:roots-proof}

\emph{Parte~(1).} Pelo \cref{cor:four}, módulo $\tfrac i2\Z$ há quatro raízes em $\Re s>0$, com multiplicidade. Pelas
\cref{prop:nk,prop:mu,prop:real}, elas são $\mu$, $\sK$, $\sphys$ (setor A) e $\sB$ (setor B, que aparece como
$\sB+\tfrac i2$ para $L_A$), todas reais e algebricamente simples, com as localizações enunciadas; o intervalo de
$\sphys$ é $\lambda_0\pm1.0163\cdot10^{-5}$, arredondado para fora com oito dígitos. Os autovetores em $\sK$ e $\sB$ violam os vínculos
(\cref{prop:witness}); o autovetor em $\mu$ é $g_*$, isto é, $e^{\mu\tau}g_*$ é o modo de gauge $G_*$
(\cref{prop:mu}); o autovetor em $\sphys$ satisfaz os vínculos (\cref{cor:sK}). O enunciado sobre o
eixo imaginário é o \cref{cor:four}, com autovetor $\partial_\tau\omegabar$ em $s=0$.

\emph{Parte~(2).} As quatro raízes são simples, de modo que a \cref{prop:physmult} se aplica na sua forma simples:
\begin{itemize}[leftmargin=*]
\item em $\sK$ e $\sB$, $C(s)h\neq0$, de modo que a multiplicidade física é $0$;
\item em $\mu$, $h=g_*$ gera $\mathcal G_\mu$, de modo que ela é $0$, e $1$ no cenário de cone fixo;
\item em $\sphys$, $C(\sphys)h=0$ e $\sphys\not\equiv\mu\pmod{\tfrac i2\Z}$, de modo que $\mathcal G_{\sphys}=0$ e ela é $1$.
\end{itemize}
Pelo \cref{thm:correspondence}, os modos de Floquet físicos com $\Re s>0$, módulo gauge, formam um espaço de dimensão um,
com expoente $\sphys$; como $\sphys$ é simples, não há modos generalizados. No cenário de cone fixo, a
contagem é dois, e o modo adicional é $\Lie_{X_0}(\gbar,\phibar)$, que, pelo \cref{lem:cone}, é tangente às
translações do ponto singular. \qed

\section{Cálculo rigoroso e reprodutibilidade}\label{sec:computation}

Esta seção descreve em que consistem as partes da prova assistidas por computador, em que é preciso confiar e como
um leitor pode conferi-las.

\subsection{Aritmética e base de confiança}\label{sec:computation-arith}

Três tipos de cálculo entram nos certificados.
\begin{enumerate}[label=(\alph*),leftmargin=*]
\item \emph{Aritmética racional exata} (o módulo \texttt{fractions} do Python): as desigualdades finais de todo
  certificado, a saber, as condições de Perron $Nv\le\theta v$ com $\theta<1$ do
  \cref{lem:perron}, as condições de Newton--Kantorovich da \cref{sec:roots-nk}, as coberturas dos contornos e
  da região da \cref{prop:Kinj}, o polinômio de deflação $q$ e o fato F1.
\item \emph{Aritmética de bolas} com Arb~\cite{Arb}, por meio do \texttt{python-flint}: as cotas dos símbolos do
  \cref{lem:F3} e a folga entre os pontos da grade.
\item \emph{Cotas a posteriori sobre cálculos em ponto flutuante.} A álgebra linear de grande porte (dimensão $14016$ para a
  referência finita, janelas de dimensão até $11200$) é feita em dupla precisão IEEE. Os resultados em ponto flutuante
  são usados apenas como inversas aproximadas, resolvedores ou vetores de teste, e toda quantidade que entra num certificado é
  cotada rigorosamente: os resíduos são calculados com cotas de arredondamento entrada a entrada da forma padrão
  $\gamma_n\||M||X|\|$~\cite{Higham2002}; as normas de inversas são cotadas por desigualdades de Loewner verificadas,
  $t^2MM^*-I\succeq0$, cuja positividade definida é verificada pelo método de Rump~\cite{Rump2006} aplicado à
  realificação; e os fatores de acoplamento são certificados como na \cref{sec:counting-reference}.
\end{enumerate}
A base de confiança é, portanto: a correção das cotas publicadas em~\cite{RT2019}, que usamos tal como enunciadas;
a aritmética IEEE~754 com arredondamento para o mais próximo, como supõem as cotas a posteriori; o Arb; e o código Python do
programa de verificação, que é curto para cada conferência e está disponível para inspeção. Nenhuma parte da cadeia final
está formalizada num assistente de provas.

Como as entradas em ponto flutuante dependem da biblioteca BLAS e do número de threads, refazer um
certificado em outra máquina reproduz o $\theta$ racional apenas a menos de $\sim10^{-16}$, não exatamente; o
critério de sucesso de uma reprodução é $\theta<1$, com uma diferença na escala do arredondamento.

\subsection{Os certificados}\label{sec:computation-certs}

A \cref{tab:certs} lista os certificados calculados, os resultados que eles estabelecem e o seu custo aproximado. Os
cálculos foram feitos em cinco computadores comuns, com quatro a seis núcleos e $7$--$31$\,GB de memória; a cadeia
inteira leva alguns dias em três ou quatro dessas máquinas. Os maiores cálculos individuais são a decomposição de Schur
da referência finita e as cotas de Loewner do seu resolvente (horas por centro do contorno), e as
caudas com $26$ janelas densas (três a seis horas por centro).

\begin{table}[ht]
\centering
\small
\begin{tabular}{@{}p{0.29\textwidth}p{0.35\textwidth}p{0.27\textwidth}@{}}
\toprule
certificado & estabelece & custo \\
\midrule
símbolos de $B$ e $B^\sharp$ (Arb, grades $2048\times4096$ e $1024\times2048$) & \cref{lem:F3}; \cref{thm:counts}~(i) & horas, 4 processos \\
recuperação e transferência do raio & \cref{lem:Lreal,lem:recovery} & minutos \\
referência finita: matriz, Schur, resíduo & \cref{lem:schur} & horas \\
nível $Z$ em $174$ discos (Loewner, resoluções de Schur) & \cref{thm:counts}~(ii),(iii) & $\sim$9 min por ponto \\
fatores de acoplamento de posto baixo ($13$ centros) & \cref{sec:counting-reference} & $\sim$1 h por centro \\
caudas, $26$ janelas ($13$ centros) & \cref{sec:counting-reference} & 3--6 h por centro \\
janelas dentro das faixas & \cref{lem:windows} & horas \\
NK no sistema aumentado em $\mu,\sK,\sphys,\sB$ & \cref{prop:nk,prop:mu} & 3--6 h cada \\
testemunhos dos vínculos & \cref{prop:witness} & minutos \\
$154$ células e Rouch\'e para $K$ & \cref{prop:Kinj} & $\sim$1 dia, 3 máquinas \\
F1 & \cref{lem:F1} & segundos \\
\bottomrule
\end{tabular}
\caption{Certificados calculados. Os nomes dos arquivos e os hashes sha256 estão listados no \cref{app:certificates}.}
\label{tab:certs}
\end{table}

\subsection{O programa de verificação}\label{sec:computation-verify}

O repositório contém um programa de verificação com dois níveis.

O \emph{nível 1} refaz, em aritmética racional exata e com o código do repositório, as desigualdades finais
dos certificados a partir dos arquivos de certificado gravados (pequenos arquivos JSON) e dos dados publicados de~\cite{RT2019}:
\begin{itemize}[leftmargin=*]
\item a cota de Perron em três níveis em cada um dos $175$ pontos dos $22$ arquivos de discos das duas faixas ($174$ discos de
  raio positivo, e uma entrada de raio zero que não é usada na cobertura);
\item as coberturas de $\partial\Omega_A$ e $\partial\Omega_B$;
\item as contagens finitas;
\item as cotas das janelas;
\item os quatro certificados de Newton--Kantorovich, os testemunhos e a identificação da raiz de gauge;
\item as $154$ células e a sua cobertura da região da \cref{prop:Kinj} fora do disco $|s-0.401|<0.05$ do
  argumento de Rouch\'e para $K$;
\item o polinômio de deflação $q$ da \cref{sec:counting-rouche};
\item as cotas de recuperação e de transferência do raio, a cota $\beta_+$ da parte do fundo e o F1.
\end{itemize}
Alguns certificados entram no nível 1 apenas pelos seus valores gravados, e só são recalculados no nível 2: as cotas dos
símbolos e das partes livres (o nível 1 recombina as faixas gravadas), e as cotas de Perron das $16$ células do
argumento de Rouch\'e para $K$.
Ele leva cerca de duas horas em quatro núcleos. Todas as conferências passam: a pior raiz de Perron nos discos é $0.999074$, com
diferenças para os valores gravados $\le7.8\cdot10^{-16}$, e, nas máquinas em que foram gerados, os dados de
Newton--Kantorovich, os testemunhos e as células são reproduzidos de forma idêntica. Em outras plataformas, as
entradas em ponto flutuante dos certificados podem diferir nos últimos bits, e o mesmo acontece com os racionais recalculados; o programa
confere então que cada certificado recalculado é válido, não passa das cotas publicadas aqui e concorda com o gravado
a menos de $10^{-9}$ relativo.

O \emph{nível 2} recalcula os certificados intermediários a partir dos dados de~\cite{RT2019}. Como amostra, recalculamos
do zero, numa máquina diferente:
\begin{itemize}[leftmargin=*]
\item a matriz de referência, idêntica bit a bit à gravada;
\item a cota do resíduo de Schur, idêntica;
\item a cauda do centro de contorno $3.0$ ($26$ janelas e o nível distante, $5.3$ horas), idêntica em todos os campos.
\end{itemize}

\subsection{Disponibilidade de dados e código}\label{sec:computation-data}

O código, os arquivos de certificado, o programa de verificação e um guia de reprodução estão disponíveis em
\url{https://github.com/edivaldo-amaral/choptuik-espectro} (licença MIT)~\cite{Repo2026}. Os arquivos de dados grandes, a saber, a matriz
de referência, os seus fatores de Schur e os fatores de acoplamento ($16$\,GiB), estão depositados no Zenodo, DOI
\href{https://doi.org/10.5281/zenodo.23222363}{10.5281/zenodo.23222363} (CC BY 4.0)~\cite{Zenodo2026}, com hashes sha256 que
também estão registrados nos arquivos de certificado. Os dados de~\cite{RT2019} não são redistribuídos; um script os baixa
do fonte de~\cite{RT2019} no arXiv e confere os seus hashes sha256.

\subsection{Conferências independentes assistidas por IA}\label{sec:computation-reviews}

Antes de serem usados, os certificados e os lemas foram conferidos por agentes de IA independentes dos que produziram o
trabalho, com as exceções listadas abaixo. Essas conferências não são revisões por especialistas humanos. O protocolo foi
desenhado para que uma conferência produza evidência verificável, e não uma opinião:
\begin{itemize}[leftmargin=*]
\item o agente conferente rodava num modelo diferente do que produziu o trabalho, com contexto limpo;
\item ele recebia as alegações enunciadas com precisão, as derivações sob revisão, o código e os dados, mas não
  as conclusões do autor nem auditorias anteriores;
\item ele tinha de entregar um artefato por alegação: uma derivação independente, uma implementação independente ou um
  cálculo direto da quantidade verdadeira para comparar com a cota;
\item antes de cada conferência era tirado um snapshot com manifesto sha256, comparado ao final dela.
\end{itemize}
Houve doze conferências desse tipo:
\begin{itemize}[leftmargin=*]
\item as células da \cref{prop:Kinj} (uma conferência anterior delas, com outra ferramenta, foi superficial e não é
  contada);
\item os certificados em $\sK$ (com o argumento de Rouch\'e para $K$), em $\mu$ e em $\sphys$;
\item os argumentos de Rouch\'e das duas faixas, incluindo o certificado em $\sB$;
\item os lemas estruturais e a cadeia da prova;
\item o lema de gauge, a reconstrução direta e a recíproca;
\item as hipóteses de modelagem;
\item o manuscrito inteiro, alegação por alegação, contra os arquivos de certificado.
\end{itemize}
Uma parte não foi conferida isoladamente: os cálculos dos símbolos do \cref{lem:F3} foram examinados apenas no seu uso,
não recalculados. Nenhuma das conferências derrubou um resultado. Todas encontraram lacunas, sobretudo de rastreabilidade ou de enunciado,
e algumas de derivação: por exemplo, uma encontrou que uma versão inicial das células usava a aproximação diádica de
$\mu$ em vez do valor verdadeiro, e outra que o erro da deflação fora da caixa finita tinha sido omitido. Todas
foram corrigidas antes de os resultados serem usados; os relatórios e artefatos estão no repositório.

\subsection{Uso de assistentes de IA}\label{sec:computation-ai}

Este trabalho foi feito com uso extensivo de assistentes de IA, entre setembro e outubro de 2026. Descrevemos o seu papel
como exige a política da revista.

\emph{Escopo.} Os assistentes escreveram a maior parte do código, fizeram a maior parte das derivações e dos cálculos, prepararam
a documentação e redigiram este manuscrito. A sua contribuição é substancial e não se limita à edição do
texto.

\emph{O papel do autor.} O autor propôs o problema, o espectro instável da solução de~\cite{RT2019}, que
os seus autores deixaram em aberto; dirigiu o projeto e escolheu entre as abordagens; rodou os cálculos nas suas próprias
máquinas e em máquinas emprestadas; conferiu os argumentos e os resultados; e assume total responsabilidade pelo conteúdo
deste artigo.

\emph{Ferramentas.} O assistente principal foi o Claude (Anthropic), usado por meio do ambiente Claude Code, com os modelos
\texttt{claude-opus-5} e, depois, \texttt{claude-opus-5-5}. O agente Codex, da OpenAI, foi usado na primeira sessão
exploratória e durante um intervalo de três dias, com os modelos \texttt{gpt-5.6-luna}, \texttt{gpt-5.6-sol},
\texttt{gpt-6-astra} e \texttt{gpt-reserve}; tudo o que foi produzido nesse intervalo foi auditado e rederivado ou
conferido antes do uso. As conferências da \cref{sec:computation-reviews} foram feitas por agentes Claude rodando
\texttt{claude-fable-5-1}, um modelo diferente do que produziu o trabalho. A pedido do autor, uma versão preliminar do
manuscrito também foi lida pelo ChatGPT (OpenAI, modelo Sol 6.1); os seus comentários foram conferidos contra os certificados e
as fontes antes de serem incorporados. Os identificadores dos modelos são os registrados nos logs das sessões; um relato datado das sessões está no arquivo \texttt{paper/AI\_USE.md} do
repositório. Esta tradução para o português também foi feita com o Claude (\texttt{claude-opus-5-5}) e revisada pelo autor.

\emph{Limites da verificação.} Os argumentos analíticos estão neste artigo; algumas estimativas de rotina aparecem em
forma de esboço, com as suas constantes conferidas por código no repositório (\cref{app:lemmas}). Toda cota calculada é
produzida por código que pode ser inspecionado e rodado de novo a partir dos arquivos de certificado e dos dados. As desigualdades finais dos
certificados são refeitas pelo programa de verificação (nível~1 da \cref{sec:computation-verify}); algumas cotas
intermediárias entram no nível~1 apenas pelos seus valores gravados, como listado na \cref{sec:computation-verify}, e só são recalculadas
no nível~2, no qual uma amostra dos certificados intermediários foi recalculada a partir dos dados
de~\cite{RT2019}. Os lemas e a maior parte dos certificados também foram
examinados pelas conferências da \cref{sec:computation-reviews}, com as exceções listadas ali. Nada disso substitui a
revisão por pares; o objetivo é que a confiança no resultado não dependa do processo que o produziu.

\section{Discussão}\label{sec:discussion}

\subsection{O que o teorema diz sobre o colapso crítico}

No quadro padrão do colapso crítico~\cite{GundlachMartinGarcia2007}, dados iniciais perto do limiar de
formação de buraco negro evoluem por muito tempo perto da solução crítica, e o desvio é dominado pelo seu
único modo que cresce, $e^{\lambda T}$. Se os dados iniciais dependem de um parâmetro $p$, a massa do buraco negro então
escala como $(p-p_*)^\gamma$ com $\gamma=1/\lambda$, modulada por uma oscilação periódica no caso discretamente
autossimilar. O Teorema Principal fornece o ingrediente espectral desse quadro para o colapso esfericamente simétrico do
campo escalar: há exatamente um modo instável, ele é simples, e $\lambda\in[2.674040,\,2.674115]$. Transformar o
quadro num teorema exige outros ingredientes, entre eles os seguintes; a questão de regularidade da
\cref{sec:discussion-regularity} e os modos neutros da \cref{rem:neutral} também teriam de ser tratados.
\begin{enumerate}[label=(\alph*),leftmargin=*]
\item \emph{Estabilidade linear no sentido da evolução.} A ausência de outras raízes de $L$ em $\Re s>0$ não dá, por
  si só, o decaimento do fluxo linearizado num complemento dos modos instável e de gauge: o operador não é
  autoadjunto, e são necessárias cotas do resolvente ao longo de retas verticais, ou um enunciado de completude para os modos de
  Floquet, e não apenas a localização dos autovalores. Reiterer deu um arcabouço abstrato para a estabilidade de
  codimensão finita de equações hiperbólicas simétricas periódicas no tempo, com resolvente compacto suficientemente à direita no plano
  complexo, motivado pelos espaços-tempos discretamente autossimilares~\cite{Reiterer2015}; se as suas hipóteses valem para a
  presente formulação não é tratado aqui.
\item \emph{Estabilidade não linear:} uma variedade estável de codimensão um da solução crítica.
\item \emph{Aspectos globais.} A solução de~\cite{RT2019} é construída numa vizinhança do cone de luz passado
  da singularidade; o quadro do limiar diz respeito a dados assintoticamente planos.
\end{enumerate}
A concordância numérica de $1/\lambda$ com o expoente medido $\gamma\approx0.374$ é, portanto, uma conferência de
consistência, e não uma consequência.

\subsection{Perturbações não esféricas}

Reiterer e Trubowitz descrevem a evidência numérica sobre perturbações gerais como inconclusiva~\cite{RT2019},
citando a análise linear de Mart\'in-Garc\'ia e Gundlach~\cite{MartinGarciaGundlach1999}, que encontraram decaimento de todas
as perturbações não esféricas, e as evoluções não lineares axissimétricas
de~\cite{ChoptuikHirschmannLieblingPretorius2003}. O setor delicado é o setor polar (par) com $\ell=2$,
cujo modo menos amortecido foi encontrado em $\kappa\Delta\approx-0.07$, perto do eixo imaginário.

Como primeiro passo para estender o Teorema Principal, escrevemos as equações de
Gerlach--Sengupta~\cite{GerlachSengupta1979} sobre o fundo de~\cite{RT2019}, nas variáveis $\omega$ e $\mu$, e
as resolvemos numericamente por colocação espectral, mantendo o cone de luz dentro do domínio e filtrando soluções
espúrias pela regularidade e pelos vínculos. Esse cálculo \emph{não} é rigoroso e não faz parte do
teorema. Nos dois setores, para $\ell=2$ e $\ell=3$, ele é consistente com os valores
de~\cite{MartinGarciaGundlach1999} dentro da precisão do cálculo por diferenças finitas deles; por exemplo,
$\kappa\Delta\approx-2.320$ contra o $-2.30$ deles para o modo axial com $\ell=2$. Para o modo polar com $\ell=2$,
ele dá
\[
  s\approx-0.004764\pm0.147345\,i,\qquad \kappa\Delta=4\pi\Re s\approx-0.0599,
\]
contra o $-0.07$ deles, e estável sob refinamento da nossa discretização até cerca de $2\cdot10^{-5}$, o que é
evidência numérica e não uma cota certificada: o modo é amortecido, mas a uma distância de apenas $0.005$ do
eixo. Uma contagem rigorosa para $\ell\ge1$ parece viável com os métodos deste artigo, usando uma deflação e uma
localização rigorosa por Newton--Kantorovich para esse modo, a estrutura dos vínculos no certificado e uma cota uniforme para
$\ell$ grande. Junto com o Teorema Principal, ela daria exatamente um modo instável para perturbações lineares gerais,
no sentido espectral.

\subsection{Regularidade}\label{sec:discussion-regularity}

A classe de modos físicos da \cref{def:physical-mode} exige analiticidade em $\xi$ através do cone de luz e
no eixo, com uma vizinhança complexa que pode ser arbitrariamente fina. Se todo modo de Floquet com
$\Re s>0$ que é apenas suave em $\xi$ é analítico é uma questão de regularidade em aberto para um sistema hiperbólico
periódico em $\tau$. A nossa recuperação do raio (\cref{lem:recovery}) funciona para todo peso $\kappa_2'>1$, mas a sua cota na
parte livre degenera como $1/(\kappa_2'-1)$, de modo que ela não alcança $\kappa_2'=1$.

\subsection{Sobre o método}

O que tornou a contagem viável em computadores comuns foi uma combinação de escolhas de formulação.
\begin{itemize}[leftmargin=*]
\item \emph{Uma referência finita em vez de um majorante do exterior.} Uma abordagem anterior cotava separadamente o resolvente
  da seção finita e o operador fora dela; ela exigia truncamentos e constantes fora de alcance
  sem computação de alto desempenho. Comparar $L$ com uma referência bloco-diagonal cujos zeros são os de uma
  matriz finita exige apenas cotas ao longo do contorno.
\item \emph{Um espaço no qual o fundo é pontual.} No espaço do toro, a imagem numérica do fundo
  se reduz ao símbolo, o que deu a cota de exclusão $2.93$, contra $414$ com majorantes bloco a bloco.
\item \emph{O critério de Perron} para matrizes de normas de blocos, em vez da norma espectral de uma tal matriz.
\item \emph{A deflação das raízes de gauge,} que estão no contorno ou perto dele.
\item \emph{Os vínculos por meio de $K$.} Num diagnóstico em ponto flutuante no truncamento $12\times36$, o máximo
  sobre a fronteira de $[1/8,1]\times[-1/4,1/4]$ da norma ponderada de $J(\,\cdot\,+s)^{-1}$ foi cerca de $1800$ para
  $L$ e $14$ para $K$, cada um no espaço em que o seu certificado é formulado. A classificação física é, portanto, obtida
  provando que $K$ é injetivo, e não cotando os vínculos nos espaços de raízes de $L$.
\end{itemize}
Várias outras rotas foram tentadas e encerradas, com cotas inferiores medidas para o seu custo. Elas estão documentadas no
repositório, para leitores que queiram tentar cálculos semelhantes.


\section*{Agradecimentos}

O autor agradece aos seus filhos, Jo\~ao Rondon e Andr\'e Leonam, por cederem o tempo ocioso dos seus
computadores, e ao Departamento de Engenharia Elétrica e Automação (DEEA) do Instituto Federal de Mato
Grosso (IFMT) pela cessão de dois computadores.

\section*{Declarações}

\emph{Financiamento.} Este trabalho não recebeu financiamento específico.

\emph{Conflito de interesses.} O autor declara não ter conflito de interesses.

\emph{Contribuições do autor.} Este é um trabalho de autor único. O autor concebeu e dirigiu o projeto e é
responsável por todo o seu conteúdo; o papel dos assistentes de IA está descrito na \cref{sec:computation-ai}.

\emph{Disponibilidade de dados e código.} O código, os arquivos de certificado e o programa de verificação estão
disponíveis no repositório~\cite{Repo2026}; os arquivos de dados grandes estão depositados no Zenodo~\cite{Zenodo2026}
(\cref{sec:computation-data}).

\emph{Uso de ferramentas de IA.} Assistentes de IA foram usados extensivamente neste trabalho; o seu papel, os modelos usados e
os procedimentos de verificação estão descritos nas \cref{sec:computation-ai,sec:computation-reviews}.


\appendix
\crefalias{section}{appendix}\crefalias{subsection}{subappendix}% amsart: sem isto o cleveref chama os apendices de secoes
\section{Provas dos lemas estruturais}\label{app:lemmas}

\subsection{Setores, deslocamento e conjugação (\cref{lem:sectorB})}\label{app:sectors}

Seja $(Uh)_{mn}=h_{m-1,n}$, isto é, $Uh=e^{i\tau/2}h$. Em $\mathbb T(a,b)$ os pesos $a^{2m}$ dão
$\|Uh\|=a\|h\|$, e em $Y(a,b)$ temos $a^{-1}\|h\|\le\|Uh\|\le a\|h\|$; assim, $U$ é um isomorfismo de cada espaço,
uma isometria a menos de um fator constante. Ele comuta com a convolução por um campo (uma soma em $m$), com a reflexão
$P$, com a multiplicação por $\xi$ e com $R_\xi$, que agem apenas em $n$ ou são invariantes por translação em
$m$; e $U^{-1}\partial_\tau U=\partial_\tau+\tfrac i2$. A parte livre $J$ satisfaz, portanto,
$U^{-1}(J+s)U=J+s+\tfrac i2$, $U$ leva o domínio de $J$ sobre si mesmo, e $U^{-1}L(s)U=L(s+\tfrac i2)$,
$U^{-1}C(s)U=C(s+\tfrac i2)$. Como $U$ muda a paridade de $m$ em toda componente, ele troca $X_A$ e $X_B$.
A conjugação por $U$ leva núcleos e cadeias de Jordan em núcleos e cadeias de Jordan, e os que satisfazem os vínculos em
outros que os satisfazem. Isso prova o \cref{lem:sectorB}.

\emph{Conjugação.} Seja $(\mathcal Ch)_{mn}=\overline{h_{-m,n}}$, que corresponde à conjugação complexa da
função $h(\tau,\xi)$ para $\tau,\xi$ reais (lembre que $h_{m,-n}=h_{mn}$). O fundo é real, de modo que
$\mathcal C L(s)\mathcal C=L(\bar s)$, e $\mathcal C$ preserva os dois setores. Em $Y(a,b)$, $\mathcal C$ é uma
isometria, de modo que as raízes de $L_A$ e de $L_B$ em $Y(a,b)$, com as suas multiplicidades, são invariantes por
$s\mapsto\bar s$. (No espaço do toro, $\mathcal C$ não é uma isometria, já que o peso $a^{2m}$ não é par em $m$;
é por isso que o argumento de realidade da \cref{prop:real} passa por $\Yp$ e pelo \cref{lem:Lreal}.) Para a faixa B,
$\sigma(L_A)\cap\{\tfrac14<\Im s<\tfrac34\}=\sigma(L_B)+\tfrac i2$, e a invariância de $\sigma(L_B)$ sob
conjugação vira invariância por $s\mapsto\overline{s-\tfrac i2}+\tfrac i2=\bar s+i$.

\subsection{A propriedade de Fredholm (\cref{prop:fredholm})}\label{app:fredholm}

\emph{(i) A parte livre.} $J+z$ age em cada modo de Fourier $m$ e em cada componente separadamente, como um operador
triangular superior em $n$ com diagonal $\mu(n+1)+z+im/2$ e estrutura fora da diagonal de primeira ordem; as inversas
explícitas de~\cite[\S3]{RT2019} são limitadas em $Y(a,b)$, com uma cota da forma $C/(\mu+z)$, e no espaço do toro são
limitadas modo a modo por formas fechadas (no repositório). A norma de $Q(I-G)$, onde $G$ é a projeção
numa caixa $\{|m|<M,\,n<N\}$, tende a zero quando $M,N\to\infty$; por exemplo, com $r=1/b$,
\[
  \|Q(I-G)\|\le\max\Big(\frac{3(1+2\mu_{\max})}{M(1-r)},\ \frac{3(1+4/(1-r))}{2\mu_{\min}(N+1)}\Big)
\]
em $Y(1,b)$, onde $\mu_{\min}\le\mu\le\mu_{\max}$. Portanto $Q$ é limite em norma de operadores de posto finito e é compacto.
O domínio maximal $\Dom_{\max}=\{u:Ju\in Y\}$, com $Ju$ calculado termo a termo, coincide com o fecho
$\Dom_{\min}$ dos polinômios trigonométricos na norma do gráfico: para $z>0$, $Q(z)$ leva $Y$ em $\Dom_{\min}$ e
inverte $J+z$ ali; para $u\in\Dom_{\max}$, $w=u-Q(z)(J+z)u$ satisfaz $(J+z)w=0$ como função analítica e, em
cada modo de Fourier, as soluções homogêneas são múltiplos de $(1+\xi)^{-1-(z+im/2)/\mu}$ (componentes com
$(1+\xi)\partial_\xi$) ou de $\xi^{-1-(z+im/2)/\mu}$ (a componente com $\xi\partial_\xi$), que são singulares em
$\xi=-1$ ou em $\xi=0$, dentro da elipse de Bernstein de parâmetro $b>1$. Assim $w=0$.

\emph{(ii) O fundo.} $B=\mu S\Gamma_1+2S\Xi\Gamma_2(\omegabar,\cdot)$ é uma soma de reflexões, da divisão regularizada
$R_\xi$, da multiplicação por $\xi$ e da convolução com os coeficientes de $\omegabar$. Estes estão em
$\ell^1$ com pesos $(65/64)^{|m|}(5/4)^{|n|}$, de modo que a convolução é limitada em todo $Y(a,b)$ e $\mathbb T(a,b)$
com $a\le65/64$, $b\le5/4$; $R_\xi$ é limitado para $b>1$~\cite[\S4]{RT2019}.

\emph{(iii)} $L(s)Q=I+(B+s-z_0)Q$ é a identidade mais um operador compacto, analítico (afim) em $s$, logo de Fredholm de
índice zero. Para $s$ real grande, $L(s)=(J+s)\big(I+(J+s)^{-1}B\big)$ é invertível em $Y(a,b)$ por uma série de Neumann, já que
$\|(J+s)^{-1}\|\le C/(\mu+s)$; no espaço do toro, $L(s)$ é invertível para $s>2.9261$ pelo \cref{thm:counts}~(i). Pelo teorema
de Fredholm analítico, as raízes são isoladas e de multiplicidade algébrica finita. O mesmo
argumento se aplica a $K(s)=J^\sharp+B^\sharp+s$ nos espaços $\sharp$.

\subsection{Recuperação do raio (\cref{lem:Lreal,lem:recovery})}\label{app:recovery}

Sejam $Y_s\subset Y_w$ uma realização forte e uma fraca (para o \cref{lem:Lreal}, $\Yp\subset\mathbb T(65/64,5/4)$;
para o \cref{lem:recovery}, $\Yp\subset Y(\kappa_1',\kappa_2')$), com o mesmo operador $Q=J^{-1}$, que preserva
caixas. Escreva $H(s)=L(s)Q=I+(B+s)Q$, seja $G$ a projeção numa caixa finita e $T=I-G$, e suponha
\[
  \|T(H(s)-I)T\|_{Y_s}<1,\qquad \|T(H(s)-I)T\|_{Y_w}<1 .
\]
\emph{Núcleos.} Seja $H(s)x=0$ com $x\in Y_w$. Então $Gx$ é um vetor finito, logo está em $Y_s$, e
$THT\cdot Tx=-THGx$, cujo lado direito está em $Y_s$ porque $H$ é limitado em $Y_s$. Em $Y_s$, o operador $THT$ na
imagem de $T$ é invertível por uma série de Neumann, de modo que existe $y\in Y_s$ com $THT\,y=-THGx$. A mesma equação vale
em $Y_w$, onde a sua solução é única; assim $Tx=y\in Y_s$, e $h=Qx$ está no domínio em $Y_s$. A inclusão
recíproca é a inclusão dos espaços. \emph{Cadeias.} Se $L(s)h_j=-h_{j-1}$, com $h_{j-1}$ já sabidamente no
domínio forte, o mesmo argumento se aplica a $x_j=Jh_j$ com o lado direito forte adicional $-Th_{j-1}$.

\emph{Constantes do \cref{lem:Lreal}.} Caixa $G=\{|m|<1024,\,n<4096\}$, $|s|\le3.1$. No espaço do toro,
$\|T(H-I)T\|\le(\|B\|+|s|)\|QT\|\le(3.9642+3.1)\cdot0.0525<0.371$, com $\|B\|$ vindo do símbolo (\cref{lem:F3}) e
$\|QT\|\le0.0525$ vindo de formas fechadas: a cauda $n\ge4096$ nos modos $|m|<400$, e a cota
$(1+2c_w/(\kappa_2-1))/(|m|/2)$ para o modo inteiro quando $|m|\ge400$. Em $\Yp$, com pesos de componentes
$\eta=(916,4096,243,1233)$ e a cota $\beta_+\le5191/500$ da parte do fundo, calculada em aritmética
exata a partir dos dados e da bola publicada de~\cite{RT2019} (\cref{app:certificates}), a mesma quantidade é
$\le0.2666$. Os pesos satisfazem $\eta_i\ge243$, de modo que $\|x\|_{\mathbb T}\le\|x\|_{\Yp}/243$.

\emph{Constantes do \cref{lem:recovery}.} No espaço fraco $Y(\kappa_1',\kappa_2')$, a parte do fundo é cotada por
$\beta(\kappa_2')=\beta_+\,D(\kappa_2')/D(5/4)$ com $D(k)=2k/(k^2-1)$, fator que vem da cota de $R_\xi$; o
peso de Fourier $\kappa_1'\in[1,65/64]$ não aumenta as normas de convolução, e $\|QT\|$ não depende dele,
porque $Q$ preserva $m$. Como $\|QT\|\to0$ para todo $\kappa_2'>1$, para todos $S$ e $\kappa_2'$ existe uma caixa com
$(\beta(\kappa_2')+S)\|QT\|<1$; caixas explícitas são calculadas, por exemplo $2^{17}\times2^{20}$ para $S=200$ e
$\kappa_2'=1+2^{-6}$. A verificação racional está no repositório (\cref{app:certificates}).

\subsection{Propagação dos vínculos (\cref{lem:KCML})}\label{app:KCML}

Sejam $O=\tfrac12(1-P)$, $E=\tfrac12(1+P)$, $X$ a multiplicação por $\xi$ e, para um resíduo $q=\Omega^+(\mu,w)$ com
componentes $(q_1,q_2,q_{2\sharp},q_3,q_4)$, ponha $F=Sq$, $c=S^\sharp q=(Oq_1,-Pq_{2\sharp})$ e $e=-Pq$. Para resíduos
do sistema, as componentes selecionadas são divisíveis por $\xi$ (os termos principal e quadrático carregam um fator
$\xi$, e as componentes selecionadas de $\Gamma_1$ são ímpares), de modo que
\[
  e=Ic+RF,\qquad Ic=(c_1,0,c_2,0,0),\qquad Rf=(XPf_1,XPf_2,0,XPf_3,XPf_4),
\]
usando $PX=-XP$ e $Pc_1=-c_1$. Reiterer e Trubowitz provam a identidade diferencial da \cref{eq:sharp} numa forma
\emph{fora da solução}~\cite[\S2]{RT2019}: com $J_s=\partial_\tau+s+\mu((1+\xi)\partial_\xi+1)$ e $\mathcal G(w,e)$
a aplicação bilinear da identidade deles,
\[
  T_s(w)e:=X\big(OJ_se_1,\ J_se_{2\sharp}-PJ_se_2\big)+\mu\big(Ee_1,\ e_1+e_2+Pe_2-2Pe_3\big)+X\mathcal G(w,e)
\]
satisfaz $T_0(w)e(w)=0$ para todo $w$ admissível, seja ou não $e(w)=0$. A única dependência de $T$ em $w$ é
por meio de $\mathcal G$, de modo que linearizar em $w$ na direção $e^{s\tau}h$ dá $T_s(w)e'+X\mathcal G(h,e(w))=0$,
onde $e'$ é o resíduo linearizado. Com $F'=L(s)h$, $c'=C(s)h$, $e'=Ic'+RF'$, e
\[
  K(s)c:=R_\xi T_s(w)Ic,\qquad M(s)f:=-R_\xi T_s(w)Rf,
\]
e $R_\xi X=I$, obtemos $K(s)C(s)h=M(s)L(s)h-R_\xi X\mathcal G(h,e(w))$. No fundo exato, $e(\omegabar)=0$,
o que é~\eqref{eq:KCML}. As formas explícitas de $K$ e $M$ estão no \cref{app:formulas}; $K$ tem a parte principal
$\operatorname{diag}(\partial_\tau+s+\mu(\xi\partial_\xi+1),\,J_s)$ de~\eqref{eq:sharp}.

\emph{Domínios.} $M(s)$ contém derivadas, de modo que a identidade vale como identidade de operadores fechados com perda de
raio. Para $1<a'<a$, $1<b'<b$, as derivadas são limitadas de $Y(a,b)$ em $Y(a',b')$:
\[
  \|\partial_\tau\|\le\frac{q_a}{2(1-q_a)^2},\qquad \|\partial_\xi\|\le\frac{2q_b}{(b'-1)(1-q_b)^2},\qquad
  q_a=a'/a,\ q_b=b'/b,
\]
que para $(a,b)=(65/64,5/4)$, $(a',b')=(129/128,9/8)$ valem no máximo $8385$ e $1440$. Portanto $C(s)$ e $M(s)$ são limitados
do domínio de $L$ em $Y(a,b)$ em $Y^\sharp(a',b')$; para $h$ nesse domínio,
$J^\sharp C(s)h=M(s)L(s)h-(B^\sharp+s)C(s)h\in Y^\sharp(a',b')$, de modo que $C(s)h$ está no domínio maximal, e portanto no fechado,
de $K$ ali, e~\eqref{eq:KCML} vale. A inclusão $\ell^1(129/128,9/8)\subset\mathbb T^\sharp(129/128,9/8)$ tem norma
$\le1$, já que $\kappa_1^m\le\kappa_1^{|m|}$ e $\sqrt{\kappa_2^{2j}+\kappa_2^{-2j}}\le2\kappa_2^j$; o argumento $\Dom_{\min}=\Dom_{\max}$
do \cref{app:fredholm} também vale no espaço do toro. Assim $C(s)h$ está no domínio no qual os certificados da
\cref{prop:Kinj} afirmam a injetividade.

\subsection{O modo de gauge (\cref{eq:gaugemode})}\label{app:gauge}

Sejam $X_0=e^{\mu\tau}Z$, $Z=\mu^{-1}\partial_\tau+\xi\partial_\xi$, $W=\mu+\xi^2Q$ e $F=\zeta+\log W$. O bloco
radial da métrica~\eqref{eq:metric} é proporcional a $e^{-2\zeta}(u_++u_-)^{-2}du_-du_+$, e $X_0$ translada
$u_\pm$; a derivada de Lie dá $\delta\zeta=e^{\mu\tau}(Z\zeta-1)$. O raio de área $e^{-\zeta}\xi/W$ é um escalar,
de modo que comparar a sua variação com a sua derivada ao longo de $X_0$ dá $\delta W=e^{\mu\tau}ZW$, logo
\[
  \delta Q=e^{\mu\tau}(ZQ+2Q),\qquad \delta F=e^{\mu\tau}(ZF-1),\qquad \delta\phi=e^{\mu\tau}Z\phi .
\]
Um cálculo direto dá $[D^\sigma,X_0]=e^{\mu\tau}D^\sigma$. Substituindo em~\eqref{eq:omega}, toda componente
tem o mesmo perfil: $\delta\omega_i=e^{\mu\tau}(Z+1)\omega_i$ para $i=2,3,4$, e também para $\omega_1=2\xi Q+2\mu\partial_\xi F$,
em que os termos extras somam $\omega_1$. Mais geralmente, para todo $w$ admissível,
\[
  D\Omega(w)\big[e^{\mu\tau}(Z+1)w\big]=e^{\mu\tau}(Z+1)\,\Omega(w),
\]
pelo comutador, pela homogeneidade quadrática de $\Gamma_2$ e pelo fator $\xi$ nos termos diferenciais e
quadráticos. No fundo exato, o lado direito se anula, o que dá $L(\mu)g_*=0$ e $C(\mu)g_*=0$.

\section{Fórmulas explícitas}\label{app:formulas}

\subsection{Pesos}
A prova de existência de~\cite{RT2019} usa $(\kappa_1,\kappa_2)=(65/64,5/4)$. O nosso espaço forte é $\Yp=Y(65/64,5/4)$
com pesos de componentes $\eta=(916,4096,243,1233)$; os certificados $\sharp$ usam $(129/128,9/8)$. O semieixo real
da elipse de Bernstein de parâmetro $5/4$ é $\tfrac12(\tfrac54+\tfrac45)=\tfrac{41}{40}$, o raio do
domínio $\Mhat$.

\subsection{A aplicação de vínculos}
Com a codificação $h_i^-=h_i$, $h_i^+=-Ph_i$ ($i=2,3,4$), $D^+=\partial_\tau+\mu(1+\xi)\partial_\xi$, e as
fórmulas~\eqref{eq:omega} e~\cite[\S2]{RT2019} para $\Omega_1^+$ e $\Omega_{2\sharp}^+$,
\[
  C(s)h=\Big(\tfrac12(1-P)\,\delta\Omega_1^+[h],\ -P\,\delta\Omega_{2\sharp}^+[h]\Big),
\]
onde
\begin{align*}
  \delta\Omega_1^+[h]&=\xi(D^++s+\mu)h_1+2\mu h_1+\tfrac14\xi\,\delta\mathcal Q_1[h],\\
  \delta\Omega_{2\sharp}^+[h]&=\xi(D^++s+\mu)h_2^++2\mu(h_2^+-h_3^+)
    +\xi\,\delta\big(-(\omega_2^+)^2+2\omega_3^+\omega_2^+-(\omega_4^+)^2\big)[h],
\end{align*}
$\delta(\cdot)[h]$ denota a derivada em $\omegabar$ na direção $h$, e
$\mathcal Q_1=\omega_1^2-2\omega_2^-\omega_1-6\omega_2^+\omega_1+4\omega_3^+\omega_1+(\omega_2^-)^2-2\omega_2^+\omega_2^-
-4\omega_3^+\omega_2^-+(\omega_2^+)^2+4\omega_3^+\omega_2^+-4(\omega_4^+)^2$. Em particular, $C(s)=C_0+sC_1$ com
\[
  C_1h=\big(0,\ -\xi h_2\big),
\]
já que $\xi h_1$ é par e $P\xi P=-\xi$.

\subsection{Os operadores $K$ e $M$}
Com $J_s=\partial_\tau+s+\mu((1+\xi)\partial_\xi+1)$, $X$ a multiplicação por $\xi$,
$\mathcal G_I=\mathcal G(\omegabar,Ic)$ e $\mathcal G_R=\mathcal G(\omegabar,Rf)$ (\cref{app:KCML}), e usando que
$c_1$ e $f_1$ são ímpares:
\begin{align*}
  (K(s)c)_1&=\big(\partial_\tau+s+\mu(\xi\partial_\xi+1)\big)c_1+(\mathcal G_I)_1,\\
  (K(s)c)_2&=J_sc_2+\mu R_\xi c_1+(\mathcal G_I)_2,\\
  (M(s)f)_1&=-\mu(\xi\partial_\xi+2)Pf_1-(\mathcal G_R)_1,\\
  (M(s)f)_2&=PJ_s(XPf_2)-\mu(Pf_1+Pf_2-f_2+2f_3)-(\mathcal G_R)_2 .
\end{align*}
$\mathcal G_I$ é o termo bilinear $\Gamma_2^\sharp(\omegabar,\cdot)$ de~\eqref{eq:sharp}. A primeira componente de $M$ não
contém $s$, e $M_1f=(0,-\xi f_2)=C_1f$, de acordo com a relação $M_1=C_1$ da \cref{sec:physical-count}.

\section{Tabela de certificados}\label{app:certificates}

O repositório contém, para cada passo da prova, a lista dos arquivos dos quais o certificado correspondente
depende, com os seus tamanhos e hashes sha256 (\texttt{build/reproducao/inventario.json}, apresentada em
\texttt{docs/REPRODUCAO\_MAPA.md}). Os arquivos grandes estão no Zenodo (\cref{sec:computation-data}), e os seus hashes
também estão em \texttt{build/reproducao/pacote\_dados.json}. A tabela abaixo dá, para cada alegação, a conferência do
programa de verificação que a refaz (\cref{sec:computation-verify}), o número de arquivos e um digest: os
primeiros $16$ dígitos hexadecimais do sha256 das linhas ordenadas ``\texttt{sha256}~~caminho'' desses arquivos, calculado por
\texttt{scripts/digest\_certificados.py}.

\begin{table}[ht]
\centering
\small
% gerado por scripts/digest_certificados.py; nao editar a mao
\begin{tabular}{@{}p{0.33\textwidth}p{0.37\textwidth}rl@{}}
\toprule
alegação & conferência do nível 1 & arquivos & digest \\
\midrule
dados de~\cite{RT2019} (\cref{sec:setting-background}) & script de download & 3 & \texttt{75d22969ac3c9d81} \\
\cref{lem:Lreal} & \texttt{L\_real\_toro}, \texttt{radius\_transfer}, \texttt{beta\_plus} & 4 & \texttt{ab394b2cb7a9cca2} \\
\cref{lem:F3}, \cref{thm:counts}~(i) & \texttt{simbolo\_L} & 10 & \texttt{d4a922dbf5a78787} \\
\cref{thm:counts}, faixa A & \texttt{discos\_A}, \texttt{cobertura\_A}, \texttt{contagem\_A}, \texttt{janelas\_A}, \texttt{deflacao\_q} & 58 & \texttt{c8a9c24901eb3993} \\
\cref{thm:counts}, faixa B; \cref{prop:nk,prop:witness} em $\sB$ & \texttt{discos\_B}, \texttt{cobertura\_B}, \texttt{contagem\_B}, \texttt{janelas\_B}, \texttt{nk\_raizB} & 86 & \texttt{a62dc9d0d789e2fc} \\
\cref{prop:mu} & \texttt{nk\_s4}, \texttt{identificacao\_s4} & 17 & \texttt{891bdaa6bcbfb3c5} \\
\cref{prop:nk,prop:witness} em $\sK$ & \texttt{nk\_s3b}, \texttt{testemunho\_s3b} & 14 & \texttt{8d44ed8d74426b3a} \\
\cref{prop:Kinj} & \texttt{s3a}, \texttt{rouche\_K}, \texttt{simbolo\_K} & 14 & \texttt{02a66ff6aff02482} \\
\cref{prop:nk} em $\sphys$ & \texttt{nk\_0733} & 19 & \texttt{31458d42757f59dd} \\
\cref{lem:F1} & \texttt{F1} & 1 & \texttt{e4a65f7b6165c1f2} \\
\cref{lem:recovery} & \texttt{R5} & 1 & \texttt{6b2d83307787fd97} \\
\bottomrule
\end{tabular}

\caption{Certificados da prova, por alegação. O digest identifica o conjunto de arquivos do qual cada alegação depende.}
\label{tab:digests}
\end{table}


\bibliographystyle{amsplain}
\bibliography{refs}

\end{document}
